% Figure: surface-tension jumps and accelerations near a level transition
% (Chapter 8). Cells, liquid polygons, PLIC segments, the faces with a
% non-zero integrated jump and the accelerations R J of both formulations are
% read from the code: generated/forces.tex (forces.jl, force_data). Labels
% hand-written.
\begin{figure}
  \centering
  \def\fxfnone{%
    \def\fxfcell##1##2##3##4##5{}\def\fxfliquid##1{}\def\fxfplic##1##2##3##4{}%
    \def\fxfjump##1##2##3##4##5##6##7##8{}\def\fxfhanging##1##2{}}%
  % Panel #1: 1 CSF, 2 CGFM.
  \def\fxfpanel#1{%
    \begin{scope}
      \clip (2.0,4.0) rectangle (5.5,6.5);
      \fxfnone
      \def\fxfcell##1##2##3##4##5{\fill[tr level ##5] (##1,##2) rectangle (##3,##4);}%
      % Generated by docs/techref/figures/generate.jl; do not edit.
\fxfcell{2}{4}{2.5}{4.5}{0}
\fxfliquid{(2.062341,4) -- (2.5,4) -- (2.5,4.5) -- (2.097185,4.5) -- cycle}
\fxfplic{2.062341}{4}{2.097185}{4.5}
\fxfcell{2.5}{4}{3}{4.5}{0}
\fxfliquid{(2.5,4) -- (3,4) -- (3,4.5) -- (2.5,4.5) -- cycle}
\fxfcell{3}{4}{3.5}{4.5}{0}
\fxfliquid{(3,4) -- (3.5,4) -- (3.5,4.5) -- (3,4.5) -- cycle}
\fxfcell{3.5}{4}{4}{4.5}{0}
\fxfliquid{(3.5,4) -- (4,4) -- (4,4.5) -- (3.5,4.5) -- cycle}
\fxfcell{2}{4.5}{2.5}{5}{0}
\fxfliquid{(2.092907,4.5) -- (2.5,4.5) -- (2.5,5) -- (2.251465,5) -- cycle}
\fxfplic{2.092907}{4.5}{2.251465}{5}
\fxfcell{2.5}{4.5}{3}{5}{0}
\fxfliquid{(2.5,4.5) -- (3,4.5) -- (3,5) -- (2.5,5) -- cycle}
\fxfcell{3}{4.5}{3.5}{5}{0}
\fxfliquid{(3,4.5) -- (3.5,4.5) -- (3.5,5) -- (3,5) -- cycle}
\fxfcell{3.5}{4.5}{4}{5}{0}
\fxfliquid{(3.5,4.5) -- (4,4.5) -- (4,5) -- (3.5,5) -- cycle}
\fxfcell{2}{5}{2.5}{5.5}{0}
\fxfliquid{(2.235883,5) -- (2.5,5) -- (2.5,5.410432) -- cycle}
\fxfplic{2.235883}{5}{2.5}{5.410432}
\fxfcell{2.5}{5}{3}{5.5}{0}
\fxfliquid{(2.5,5) -- (3,5) -- (3,5.5) -- (2.558302,5.5) -- (2.5,5.4094) -- cycle}
\fxfplic{2.558302}{5.5}{2.5}{5.4094}
\fxfcell{3}{5}{3.5}{5.5}{0}
\fxfliquid{(3,5) -- (3.5,5) -- (3.5,5.5) -- (3,5.5) -- cycle}
\fxfcell{3.5}{5}{4}{5.5}{0}
\fxfliquid{(3.5,5) -- (4,5) -- (4,5.5) -- (3.5,5.5) -- cycle}
\fxfcell{2}{5.5}{2.5}{6}{0}
\fxfcell{2.5}{5.5}{3}{6}{0}
\fxfliquid{(2.5,5.5) -- (3,5.5) -- (3,5.880728) -- (2.5,5.510641) -- cycle}
\fxfplic{3}{5.880728}{2.5}{5.510641}
\fxfcell{3}{5.5}{3.5}{6}{0}
\fxfliquid{(3,5.5) -- (3.5,5.5) -- (3.5,6) -- (3.141854,6) -- (3,5.919773) -- cycle}
\fxfplic{3.141854}{6}{3}{5.919773}
\fxfcell{3.5}{5.5}{4}{6}{0}
\fxfliquid{(3.5,5.5) -- (4,5.5) -- (4,6) -- (3.5,6) -- cycle}
\fxfcell{2}{6}{2.5}{6.5}{0}
\fxfcell{2.5}{6}{3}{6.5}{0}
\fxfcell{3}{6}{3.5}{6.5}{0}
\fxfliquid{(3.143521,6) -- (3.5,6) -- (3.5,6.201611) -- cycle}
\fxfplic{3.143521}{6}{3.5}{6.201611}
\fxfcell{3.5}{6}{4}{6.5}{0}
\fxfliquid{(3.5,6) -- (4,6) -- (4,6.322031) -- (3.5,6.189898) -- cycle}
\fxfplic{4}{6.322031}{3.5}{6.189898}
\fxfcell{4}{4}{4.25}{4.25}{1}
\fxfliquid{(4,4) -- (4.25,4) -- (4.25,4.25) -- (4,4.25) -- cycle}
\fxfcell{4.25}{4}{4.5}{4.25}{1}
\fxfliquid{(4.25,4) -- (4.5,4) -- (4.5,4.25) -- (4.25,4.25) -- cycle}
\fxfcell{4.5}{4}{4.75}{4.25}{1}
\fxfliquid{(4.5,4) -- (4.75,4) -- (4.75,4.25) -- (4.5,4.25) -- cycle}
\fxfcell{4.75}{4}{5}{4.25}{1}
\fxfliquid{(4.75,4) -- (5,4) -- (5,4.25) -- (4.75,4.25) -- cycle}
\fxfcell{4}{4.25}{4.25}{4.5}{1}
\fxfliquid{(4,4.25) -- (4.25,4.25) -- (4.25,4.5) -- (4,4.5) -- cycle}
\fxfcell{4.25}{4.25}{4.5}{4.5}{1}
\fxfliquid{(4.25,4.25) -- (4.5,4.25) -- (4.5,4.5) -- (4.25,4.5) -- cycle}
\fxfcell{4.5}{4.25}{4.75}{4.5}{1}
\fxfliquid{(4.5,4.25) -- (4.75,4.25) -- (4.75,4.5) -- (4.5,4.5) -- cycle}
\fxfcell{4.75}{4.25}{5}{4.5}{1}
\fxfliquid{(4.75,4.25) -- (5,4.25) -- (5,4.5) -- (4.75,4.5) -- cycle}
\fxfcell{4}{4.5}{4.25}{4.75}{1}
\fxfliquid{(4,4.5) -- (4.25,4.5) -- (4.25,4.75) -- (4,4.75) -- cycle}
\fxfcell{4.25}{4.5}{4.5}{4.75}{1}
\fxfliquid{(4.25,4.5) -- (4.5,4.5) -- (4.5,4.75) -- (4.25,4.75) -- cycle}
\fxfcell{4.5}{4.5}{4.75}{4.75}{1}
\fxfliquid{(4.5,4.5) -- (4.75,4.5) -- (4.75,4.75) -- (4.5,4.75) -- cycle}
\fxfcell{4.75}{4.5}{5}{4.75}{1}
\fxfliquid{(4.75,4.5) -- (5,4.5) -- (5,4.75) -- (4.75,4.75) -- cycle}
\fxfcell{4}{4.75}{4.25}{5}{1}
\fxfliquid{(4,4.75) -- (4.25,4.75) -- (4.25,5) -- (4,5) -- cycle}
\fxfcell{4.25}{4.75}{4.5}{5}{1}
\fxfliquid{(4.25,4.75) -- (4.5,4.75) -- (4.5,5) -- (4.25,5) -- cycle}
\fxfcell{4.5}{4.75}{4.75}{5}{1}
\fxfliquid{(4.5,4.75) -- (4.75,4.75) -- (4.75,5) -- (4.5,5) -- cycle}
\fxfcell{4.75}{4.75}{5}{5}{1}
\fxfliquid{(4.75,4.75) -- (5,4.75) -- (5,5) -- (4.75,5) -- cycle}
\fxfcell{5}{4}{5.25}{4.25}{1}
\fxfliquid{(5,4) -- (5.25,4) -- (5.25,4.25) -- (5,4.25) -- cycle}
\fxfcell{5.25}{4}{5.5}{4.25}{1}
\fxfliquid{(5.25,4) -- (5.5,4) -- (5.5,4.25) -- (5.25,4.25) -- cycle}
\fxfcell{5}{4.25}{5.25}{4.5}{1}
\fxfliquid{(5,4.25) -- (5.25,4.25) -- (5.25,4.5) -- (5,4.5) -- cycle}
\fxfcell{5.25}{4.25}{5.5}{4.5}{1}
\fxfliquid{(5.25,4.25) -- (5.5,4.25) -- (5.5,4.5) -- (5.25,4.5) -- cycle}
\fxfcell{5}{4.5}{5.25}{4.75}{1}
\fxfliquid{(5,4.5) -- (5.25,4.5) -- (5.25,4.75) -- (5,4.75) -- cycle}
\fxfcell{5.25}{4.5}{5.5}{4.75}{1}
\fxfliquid{(5.25,4.5) -- (5.5,4.5) -- (5.5,4.75) -- (5.25,4.75) -- cycle}
\fxfcell{5}{4.75}{5.25}{5}{1}
\fxfliquid{(5,4.75) -- (5.25,4.75) -- (5.25,5) -- (5,5) -- cycle}
\fxfcell{5.25}{4.75}{5.5}{5}{1}
\fxfliquid{(5.25,4.75) -- (5.5,4.75) -- (5.5,5) -- (5.25,5) -- cycle}
\fxfcell{4}{5}{4.25}{5.25}{1}
\fxfliquid{(4,5) -- (4.25,5) -- (4.25,5.25) -- (4,5.25) -- cycle}
\fxfcell{4.25}{5}{4.5}{5.25}{1}
\fxfliquid{(4.25,5) -- (4.5,5) -- (4.5,5.25) -- (4.25,5.25) -- cycle}
\fxfcell{4.5}{5}{4.75}{5.25}{1}
\fxfliquid{(4.5,5) -- (4.75,5) -- (4.75,5.25) -- (4.5,5.25) -- cycle}
\fxfcell{4.75}{5}{5}{5.25}{1}
\fxfliquid{(4.75,5) -- (5,5) -- (5,5.25) -- (4.75,5.25) -- cycle}
\fxfcell{4}{5.25}{4.25}{5.5}{1}
\fxfliquid{(4,5.25) -- (4.25,5.25) -- (4.25,5.5) -- (4,5.5) -- cycle}
\fxfcell{4.25}{5.25}{4.5}{5.5}{1}
\fxfliquid{(4.25,5.25) -- (4.5,5.25) -- (4.5,5.5) -- (4.25,5.5) -- cycle}
\fxfcell{4.5}{5.25}{4.75}{5.5}{1}
\fxfliquid{(4.5,5.25) -- (4.75,5.25) -- (4.75,5.5) -- (4.5,5.5) -- cycle}
\fxfcell{4.75}{5.25}{5}{5.5}{1}
\fxfliquid{(4.75,5.25) -- (5,5.25) -- (5,5.5) -- (4.75,5.5) -- cycle}
\fxfcell{4}{5.5}{4.25}{5.75}{1}
\fxfliquid{(4,5.5) -- (4.25,5.5) -- (4.25,5.75) -- (4,5.75) -- cycle}
\fxfcell{4.25}{5.5}{4.5}{5.75}{1}
\fxfliquid{(4.25,5.5) -- (4.5,5.5) -- (4.5,5.75) -- (4.25,5.75) -- cycle}
\fxfcell{4.5}{5.5}{4.75}{5.75}{1}
\fxfliquid{(4.5,5.5) -- (4.75,5.5) -- (4.75,5.75) -- (4.5,5.75) -- cycle}
\fxfcell{4.75}{5.5}{5}{5.75}{1}
\fxfliquid{(4.75,5.5) -- (5,5.5) -- (5,5.75) -- (4.75,5.75) -- cycle}
\fxfcell{4}{5.75}{4.25}{6}{1}
\fxfliquid{(4,5.75) -- (4.25,5.75) -- (4.25,6) -- (4,6) -- cycle}
\fxfcell{4.25}{5.75}{4.5}{6}{1}
\fxfliquid{(4.25,5.75) -- (4.5,5.75) -- (4.5,6) -- (4.25,6) -- cycle}
\fxfcell{4.5}{5.75}{4.75}{6}{1}
\fxfliquid{(4.5,5.75) -- (4.75,5.75) -- (4.75,6) -- (4.5,6) -- cycle}
\fxfcell{4.75}{5.75}{5}{6}{1}
\fxfliquid{(4.75,5.75) -- (5,5.75) -- (5,6) -- (4.75,6) -- cycle}
\fxfcell{5}{5}{5.25}{5.25}{1}
\fxfliquid{(5,5) -- (5.25,5) -- (5.25,5.25) -- (5,5.25) -- cycle}
\fxfcell{5.25}{5}{5.5}{5.25}{1}
\fxfliquid{(5.25,5) -- (5.5,5) -- (5.5,5.25) -- (5.25,5.25) -- cycle}
\fxfcell{5}{5.25}{5.25}{5.5}{1}
\fxfliquid{(5,5.25) -- (5.25,5.25) -- (5.25,5.5) -- (5,5.5) -- cycle}
\fxfcell{5.25}{5.25}{5.5}{5.5}{1}
\fxfliquid{(5.25,5.25) -- (5.5,5.25) -- (5.5,5.5) -- (5.25,5.5) -- cycle}
\fxfcell{5}{5.5}{5.25}{5.75}{1}
\fxfliquid{(5,5.5) -- (5.25,5.5) -- (5.25,5.75) -- (5,5.75) -- cycle}
\fxfcell{5.25}{5.5}{5.5}{5.75}{1}
\fxfliquid{(5.25,5.5) -- (5.5,5.5) -- (5.5,5.75) -- (5.25,5.75) -- cycle}
\fxfcell{5}{5.75}{5.25}{6}{1}
\fxfliquid{(5,5.75) -- (5.25,5.75) -- (5.25,6) -- (5,6) -- cycle}
\fxfcell{5.25}{5.75}{5.5}{6}{1}
\fxfliquid{(5.25,5.75) -- (5.5,5.75) -- (5.5,5.980812) -- (5.465599,6) -- (5.25,6) -- cycle}
\fxfplic{5.5}{5.980812}{5.465599}{6}
\fxfcell{4}{6}{4.25}{6.25}{1}
\fxfliquid{(4,6) -- (4.25,6) -- (4.25,6.25) -- (4,6.25) -- cycle}
\fxfcell{4.25}{6}{4.5}{6.25}{1}
\fxfliquid{(4.25,6) -- (4.5,6) -- (4.5,6.25) -- (4.25,6.25) -- cycle}
\fxfcell{4.5}{6}{4.75}{6.25}{1}
\fxfliquid{(4.5,6) -- (4.75,6) -- (4.75,6.25) -- (4.5,6.25) -- cycle}
\fxfcell{4.75}{6}{5}{6.25}{1}
\fxfliquid{(4.75,6) -- (5,6) -- (5,6.219618) -- (4.887192,6.25) -- (4.75,6.25) -- cycle}
\fxfplic{5}{6.219618}{4.887192}{6.25}
\fxfcell{4}{6.25}{4.25}{6.5}{1}
\fxfliquid{(4,6.25) -- (4.25,6.25) -- (4.25,6.331586) -- (4,6.312301) -- cycle}
\fxfplic{4.25}{6.331586}{4}{6.312301}
\fxfcell{4.25}{6.25}{4.5}{6.5}{1}
\fxfliquid{(4.25,6.25) -- (4.5,6.25) -- (4.5,6.323328) -- (4.25,6.331808) -- cycle}
\fxfplic{4.5}{6.323328}{4.25}{6.331808}
\fxfcell{4.5}{6.25}{4.75}{6.5}{1}
\fxfliquid{(4.5,6.25) -- (4.75,6.25) -- (4.75,6.286416) -- (4.5,6.323551) -- cycle}
\fxfplic{4.75}{6.286416}{4.5}{6.323551}
\fxfcell{4.75}{6.25}{5}{6.5}{1}
\fxfliquid{(4.75,6.25) -- (4.887292,6.25) -- (4.75,6.286976) -- cycle}
\fxfplic{4.887292}{6.25}{4.75}{6.286976}
\fxfcell{5}{6}{5.25}{6.25}{1}
\fxfliquid{(5,6) -- (5.25,6) -- (5.25,6.12001) -- (5,6.220636) -- cycle}
\fxfplic{5.25}{6.12001}{5}{6.220636}
\fxfcell{5.25}{6}{5.5}{6.25}{1}
\fxfliquid{(5.25,6) -- (5.468716,6) -- (5.25,6.121998) -- cycle}
\fxfplic{5.468716}{6}{5.25}{6.121998}
\fxfcell{5}{6.25}{5.25}{6.5}{1}
\fxfcell{5.25}{6.25}{5.5}{6.5}{1}
\fxfjump{1}{1}{2}{4.25}{0.5}{0.134521}{0.397586}{0}
\fxfjump{1}{1}{2}{4.75}{0.5}{0.105154}{0.398146}{0}
\fxfjump{1}{1}{2}{5.25}{0.5}{0.035039}{0.398737}{0}
\fxfjump{1}{1}{2.5}{4.25}{0.5}{0.025533}{0.034506}{0}
\fxfjump{1}{1}{2.5}{4.75}{0.5}{0.055233}{0.082967}{0}
\fxfjump{1}{1}{2.5}{5.25}{0.5}{0.124868}{0.257334}{0}
\fxfjump{1}{1}{2.5}{5.75}{0.5}{0.060092}{0.380372}{0}
\fxfjump{1}{1}{3}{5.25}{0.5}{0.001707}{0.002135}{0}
\fxfjump{1}{1}{3}{5.75}{0.5}{0.090339}{0.16412}{0}
\fxfjump{1}{1}{3}{6.25}{0.5}{0.023173}{0.395913}{0}
\fxfjump{1}{1}{3.5}{5.75}{0.5}{0.003669}{0.004616}{0}
\fxfjump{1}{1}{3.5}{6.25}{0.5}{0.059127}{0.223858}{0}
\fxfjump{1}{1}{4}{5.875}{0.25}{0}{0.022437}{1}
\fxfjump{1}{1}{4}{6.125}{0.25}{0.039114}{0.101579}{0}
\fxfjump{1}{1}{4}{6.375}{0.25}{-0.017527}{-0.006382}{0}
\fxfjump{1}{1}{4.25}{6.375}{0.25}{0.001741}{0.028892}{0}
\fxfjump{1}{1}{4.5}{6.375}{0.25}{-0.007123}{-0.133311}{0}
\fxfjump{1}{1}{4.75}{6.125}{0.25}{-0.002163}{-0.010909}{0}
\fxfjump{1}{1}{4.75}{6.375}{0.25}{-0.014142}{-0.536472}{0}
\fxfjump{1}{1}{5}{6.125}{0.25}{-0.022985}{-0.138253}{0}
\fxfjump{1}{1}{5}{6.375}{0.25}{-0.003203}{-0.748703}{0}
\fxfjump{1}{1}{5.25}{5.875}{0.25}{-0.000417}{-0.00208}{0}
\fxfjump{1}{1}{5.25}{6.125}{0.25}{-0.036929}{-0.410157}{0}
\fxfjump{1}{1}{5.5}{5.875}{0.25}{-0.033923}{-0.216249}{0}
\fxfjump{1}{1}{5.5}{6.125}{0.25}{-0.016857}{-0.779232}{0}
\fxfjump{1}{2}{2.25}{4}{0.5}{0.007279}{0.011069}{0}
\fxfjump{1}{2}{2.25}{4.5}{0.5}{-0.029624}{-0.049238}{0}
\fxfjump{1}{2}{2.25}{5}{0.5}{-0.070613}{-0.201073}{0}
\fxfjump{1}{2}{2.75}{5}{0.5}{-0.001707}{-0.002135}{0}
\fxfjump{1}{2}{2.25}{5.5}{0.5}{-0.035039}{-0.398737}{0}
\fxfjump{1}{2}{2.75}{5.5}{0.5}{-0.092031}{-0.165718}{0}
\fxfjump{1}{2}{3.25}{5.5}{0.5}{-0.003669}{-0.004616}{0}
\fxfjump{1}{2}{5.375}{5.75}{0.25}{-0.000417}{-0.00208}{0}
\fxfjump{1}{2}{2.75}{6}{0.5}{-0.060092}{-0.380372}{0}
\fxfjump{1}{2}{3.25}{6}{0.5}{-0.134374}{-0.297947}{0}
\fxfjump{1}{2}{3.75}{6}{0.5}{-0.078227}{-0.119258}{0}
\fxfjump{1}{2}{4.125}{6}{0.25}{0}{-0.015671}{1}
\fxfjump{1}{2}{4.875}{6}{0.25}{-0.002163}{-0.010909}{0}
\fxfjump{1}{2}{5.125}{6}{0.25}{-0.02515}{-0.14881}{0}
\fxfjump{1}{2}{5.375}{6}{0.25}{-0.061696}{-0.507764}{0}
\fxfjump{1}{2}{4.125}{6.25}{0.25}{-0.054001}{-0.405269}{0}
\fxfjump{1}{2}{4.375}{6.25}{0.25}{-0.054385}{-0.412772}{0}
\fxfjump{1}{2}{4.625}{6.25}{0.25}{-0.061514}{-0.501399}{0}
\fxfjump{1}{2}{4.875}{6.25}{0.25}{-0.073511}{-0.721207}{0}
\fxfjump{1}{2}{5.125}{6.25}{0.25}{-0.053762}{-0.783653}{0}
\fxfjump{1}{2}{5.375}{6.25}{0.25}{-0.016857}{-0.779232}{0}
\fxfjump{1}{2}{3.25}{6.5}{0.5}{-0.023173}{-0.395913}{0}
\fxfjump{1}{2}{3.75}{6.5}{0.5}{-0.082051}{-0.408191}{0}
\fxfjump{1}{2}{4.125}{6.5}{0.25}{-0.021819}{-0.718584}{0}
\fxfjump{1}{2}{4.375}{6.5}{0.25}{-0.024465}{-0.780304}{0}
\fxfjump{1}{2}{4.625}{6.5}{0.25}{-0.017343}{-0.778333}{0}
\fxfjump{1}{2}{4.875}{6.5}{0.25}{-0.003203}{-0.748703}{0}
\fxfjump{2}{1}{2}{4.25}{0.5}{0.160054}{0.214311}{0}
\fxfjump{2}{1}{2}{4.75}{0.5}{0.160387}{0.27512}{0}
\fxfjump{2}{1}{2.5}{5.25}{0.5}{0.161615}{0.15089}{0}
\fxfjump{2}{1}{3}{5.75}{0.5}{0.154197}{0.12691}{0}
\fxfjump{2}{1}{3.5}{6.25}{0.5}{0.160588}{0.275447}{0}
\fxfjump{2}{1}{4}{6.125}{0.25}{0}{-0.035963}{1}
\fxfjump{2}{1}{4}{6.375}{0.25}{-0.078193}{-0.175936}{0}
\fxfjump{2}{1}{5.25}{6.125}{0.25}{-0.078937}{-0.397193}{0}
\fxfjump{2}{2}{2.25}{5}{0.5}{-0.160914}{-0.207587}{0}
\fxfjump{2}{2}{2.75}{5.5}{0.5}{-0.153881}{-0.125533}{0}
\fxfjump{2}{2}{3.25}{6}{0.5}{-0.161216}{-0.161947}{0}
\fxfjump{2}{2}{5.375}{6}{0.25}{-0.07897}{-0.294446}{0}
\fxfjump{2}{2}{4.125}{6.25}{0.25}{-0.075821}{-0.246759}{0}
\fxfjump{2}{2}{4.375}{6.25}{0.25}{-0.078849}{-0.271125}{0}
\fxfjump{2}{2}{4.625}{6.25}{0.25}{-0.078857}{-0.291199}{0}
\fxfjump{2}{2}{4.875}{6.25}{0.25}{-0.078877}{-0.350587}{0}
\fxfjump{2}{2}{5.125}{6.25}{0.25}{-0.078911}{-0.521108}{0}
\fxfjump{2}{2}{3.75}{6.5}{0.5}{-0.160278}{-0.33277}{0}
\fxfjump{2}{2}{4.125}{6.5}{0.25}{0}{-0.067844}{1}
\fxfhanging{4}{4.25}
\fxfhanging{4}{4.75}
\fxfhanging{4}{5.25}
\fxfhanging{4}{5.75}
\fxfhanging{4}{6.25}
%
      \fxfnone
      \def\fxfliquid##1{\fill[tr liquid] ##1;}%
      % Generated by docs/techref/figures/generate.jl; do not edit.
\fxfcell{2}{4}{2.5}{4.5}{0}
\fxfliquid{(2.062341,4) -- (2.5,4) -- (2.5,4.5) -- (2.097185,4.5) -- cycle}
\fxfplic{2.062341}{4}{2.097185}{4.5}
\fxfcell{2.5}{4}{3}{4.5}{0}
\fxfliquid{(2.5,4) -- (3,4) -- (3,4.5) -- (2.5,4.5) -- cycle}
\fxfcell{3}{4}{3.5}{4.5}{0}
\fxfliquid{(3,4) -- (3.5,4) -- (3.5,4.5) -- (3,4.5) -- cycle}
\fxfcell{3.5}{4}{4}{4.5}{0}
\fxfliquid{(3.5,4) -- (4,4) -- (4,4.5) -- (3.5,4.5) -- cycle}
\fxfcell{2}{4.5}{2.5}{5}{0}
\fxfliquid{(2.092907,4.5) -- (2.5,4.5) -- (2.5,5) -- (2.251465,5) -- cycle}
\fxfplic{2.092907}{4.5}{2.251465}{5}
\fxfcell{2.5}{4.5}{3}{5}{0}
\fxfliquid{(2.5,4.5) -- (3,4.5) -- (3,5) -- (2.5,5) -- cycle}
\fxfcell{3}{4.5}{3.5}{5}{0}
\fxfliquid{(3,4.5) -- (3.5,4.5) -- (3.5,5) -- (3,5) -- cycle}
\fxfcell{3.5}{4.5}{4}{5}{0}
\fxfliquid{(3.5,4.5) -- (4,4.5) -- (4,5) -- (3.5,5) -- cycle}
\fxfcell{2}{5}{2.5}{5.5}{0}
\fxfliquid{(2.235883,5) -- (2.5,5) -- (2.5,5.410432) -- cycle}
\fxfplic{2.235883}{5}{2.5}{5.410432}
\fxfcell{2.5}{5}{3}{5.5}{0}
\fxfliquid{(2.5,5) -- (3,5) -- (3,5.5) -- (2.558302,5.5) -- (2.5,5.4094) -- cycle}
\fxfplic{2.558302}{5.5}{2.5}{5.4094}
\fxfcell{3}{5}{3.5}{5.5}{0}
\fxfliquid{(3,5) -- (3.5,5) -- (3.5,5.5) -- (3,5.5) -- cycle}
\fxfcell{3.5}{5}{4}{5.5}{0}
\fxfliquid{(3.5,5) -- (4,5) -- (4,5.5) -- (3.5,5.5) -- cycle}
\fxfcell{2}{5.5}{2.5}{6}{0}
\fxfcell{2.5}{5.5}{3}{6}{0}
\fxfliquid{(2.5,5.5) -- (3,5.5) -- (3,5.880728) -- (2.5,5.510641) -- cycle}
\fxfplic{3}{5.880728}{2.5}{5.510641}
\fxfcell{3}{5.5}{3.5}{6}{0}
\fxfliquid{(3,5.5) -- (3.5,5.5) -- (3.5,6) -- (3.141854,6) -- (3,5.919773) -- cycle}
\fxfplic{3.141854}{6}{3}{5.919773}
\fxfcell{3.5}{5.5}{4}{6}{0}
\fxfliquid{(3.5,5.5) -- (4,5.5) -- (4,6) -- (3.5,6) -- cycle}
\fxfcell{2}{6}{2.5}{6.5}{0}
\fxfcell{2.5}{6}{3}{6.5}{0}
\fxfcell{3}{6}{3.5}{6.5}{0}
\fxfliquid{(3.143521,6) -- (3.5,6) -- (3.5,6.201611) -- cycle}
\fxfplic{3.143521}{6}{3.5}{6.201611}
\fxfcell{3.5}{6}{4}{6.5}{0}
\fxfliquid{(3.5,6) -- (4,6) -- (4,6.322031) -- (3.5,6.189898) -- cycle}
\fxfplic{4}{6.322031}{3.5}{6.189898}
\fxfcell{4}{4}{4.25}{4.25}{1}
\fxfliquid{(4,4) -- (4.25,4) -- (4.25,4.25) -- (4,4.25) -- cycle}
\fxfcell{4.25}{4}{4.5}{4.25}{1}
\fxfliquid{(4.25,4) -- (4.5,4) -- (4.5,4.25) -- (4.25,4.25) -- cycle}
\fxfcell{4.5}{4}{4.75}{4.25}{1}
\fxfliquid{(4.5,4) -- (4.75,4) -- (4.75,4.25) -- (4.5,4.25) -- cycle}
\fxfcell{4.75}{4}{5}{4.25}{1}
\fxfliquid{(4.75,4) -- (5,4) -- (5,4.25) -- (4.75,4.25) -- cycle}
\fxfcell{4}{4.25}{4.25}{4.5}{1}
\fxfliquid{(4,4.25) -- (4.25,4.25) -- (4.25,4.5) -- (4,4.5) -- cycle}
\fxfcell{4.25}{4.25}{4.5}{4.5}{1}
\fxfliquid{(4.25,4.25) -- (4.5,4.25) -- (4.5,4.5) -- (4.25,4.5) -- cycle}
\fxfcell{4.5}{4.25}{4.75}{4.5}{1}
\fxfliquid{(4.5,4.25) -- (4.75,4.25) -- (4.75,4.5) -- (4.5,4.5) -- cycle}
\fxfcell{4.75}{4.25}{5}{4.5}{1}
\fxfliquid{(4.75,4.25) -- (5,4.25) -- (5,4.5) -- (4.75,4.5) -- cycle}
\fxfcell{4}{4.5}{4.25}{4.75}{1}
\fxfliquid{(4,4.5) -- (4.25,4.5) -- (4.25,4.75) -- (4,4.75) -- cycle}
\fxfcell{4.25}{4.5}{4.5}{4.75}{1}
\fxfliquid{(4.25,4.5) -- (4.5,4.5) -- (4.5,4.75) -- (4.25,4.75) -- cycle}
\fxfcell{4.5}{4.5}{4.75}{4.75}{1}
\fxfliquid{(4.5,4.5) -- (4.75,4.5) -- (4.75,4.75) -- (4.5,4.75) -- cycle}
\fxfcell{4.75}{4.5}{5}{4.75}{1}
\fxfliquid{(4.75,4.5) -- (5,4.5) -- (5,4.75) -- (4.75,4.75) -- cycle}
\fxfcell{4}{4.75}{4.25}{5}{1}
\fxfliquid{(4,4.75) -- (4.25,4.75) -- (4.25,5) -- (4,5) -- cycle}
\fxfcell{4.25}{4.75}{4.5}{5}{1}
\fxfliquid{(4.25,4.75) -- (4.5,4.75) -- (4.5,5) -- (4.25,5) -- cycle}
\fxfcell{4.5}{4.75}{4.75}{5}{1}
\fxfliquid{(4.5,4.75) -- (4.75,4.75) -- (4.75,5) -- (4.5,5) -- cycle}
\fxfcell{4.75}{4.75}{5}{5}{1}
\fxfliquid{(4.75,4.75) -- (5,4.75) -- (5,5) -- (4.75,5) -- cycle}
\fxfcell{5}{4}{5.25}{4.25}{1}
\fxfliquid{(5,4) -- (5.25,4) -- (5.25,4.25) -- (5,4.25) -- cycle}
\fxfcell{5.25}{4}{5.5}{4.25}{1}
\fxfliquid{(5.25,4) -- (5.5,4) -- (5.5,4.25) -- (5.25,4.25) -- cycle}
\fxfcell{5}{4.25}{5.25}{4.5}{1}
\fxfliquid{(5,4.25) -- (5.25,4.25) -- (5.25,4.5) -- (5,4.5) -- cycle}
\fxfcell{5.25}{4.25}{5.5}{4.5}{1}
\fxfliquid{(5.25,4.25) -- (5.5,4.25) -- (5.5,4.5) -- (5.25,4.5) -- cycle}
\fxfcell{5}{4.5}{5.25}{4.75}{1}
\fxfliquid{(5,4.5) -- (5.25,4.5) -- (5.25,4.75) -- (5,4.75) -- cycle}
\fxfcell{5.25}{4.5}{5.5}{4.75}{1}
\fxfliquid{(5.25,4.5) -- (5.5,4.5) -- (5.5,4.75) -- (5.25,4.75) -- cycle}
\fxfcell{5}{4.75}{5.25}{5}{1}
\fxfliquid{(5,4.75) -- (5.25,4.75) -- (5.25,5) -- (5,5) -- cycle}
\fxfcell{5.25}{4.75}{5.5}{5}{1}
\fxfliquid{(5.25,4.75) -- (5.5,4.75) -- (5.5,5) -- (5.25,5) -- cycle}
\fxfcell{4}{5}{4.25}{5.25}{1}
\fxfliquid{(4,5) -- (4.25,5) -- (4.25,5.25) -- (4,5.25) -- cycle}
\fxfcell{4.25}{5}{4.5}{5.25}{1}
\fxfliquid{(4.25,5) -- (4.5,5) -- (4.5,5.25) -- (4.25,5.25) -- cycle}
\fxfcell{4.5}{5}{4.75}{5.25}{1}
\fxfliquid{(4.5,5) -- (4.75,5) -- (4.75,5.25) -- (4.5,5.25) -- cycle}
\fxfcell{4.75}{5}{5}{5.25}{1}
\fxfliquid{(4.75,5) -- (5,5) -- (5,5.25) -- (4.75,5.25) -- cycle}
\fxfcell{4}{5.25}{4.25}{5.5}{1}
\fxfliquid{(4,5.25) -- (4.25,5.25) -- (4.25,5.5) -- (4,5.5) -- cycle}
\fxfcell{4.25}{5.25}{4.5}{5.5}{1}
\fxfliquid{(4.25,5.25) -- (4.5,5.25) -- (4.5,5.5) -- (4.25,5.5) -- cycle}
\fxfcell{4.5}{5.25}{4.75}{5.5}{1}
\fxfliquid{(4.5,5.25) -- (4.75,5.25) -- (4.75,5.5) -- (4.5,5.5) -- cycle}
\fxfcell{4.75}{5.25}{5}{5.5}{1}
\fxfliquid{(4.75,5.25) -- (5,5.25) -- (5,5.5) -- (4.75,5.5) -- cycle}
\fxfcell{4}{5.5}{4.25}{5.75}{1}
\fxfliquid{(4,5.5) -- (4.25,5.5) -- (4.25,5.75) -- (4,5.75) -- cycle}
\fxfcell{4.25}{5.5}{4.5}{5.75}{1}
\fxfliquid{(4.25,5.5) -- (4.5,5.5) -- (4.5,5.75) -- (4.25,5.75) -- cycle}
\fxfcell{4.5}{5.5}{4.75}{5.75}{1}
\fxfliquid{(4.5,5.5) -- (4.75,5.5) -- (4.75,5.75) -- (4.5,5.75) -- cycle}
\fxfcell{4.75}{5.5}{5}{5.75}{1}
\fxfliquid{(4.75,5.5) -- (5,5.5) -- (5,5.75) -- (4.75,5.75) -- cycle}
\fxfcell{4}{5.75}{4.25}{6}{1}
\fxfliquid{(4,5.75) -- (4.25,5.75) -- (4.25,6) -- (4,6) -- cycle}
\fxfcell{4.25}{5.75}{4.5}{6}{1}
\fxfliquid{(4.25,5.75) -- (4.5,5.75) -- (4.5,6) -- (4.25,6) -- cycle}
\fxfcell{4.5}{5.75}{4.75}{6}{1}
\fxfliquid{(4.5,5.75) -- (4.75,5.75) -- (4.75,6) -- (4.5,6) -- cycle}
\fxfcell{4.75}{5.75}{5}{6}{1}
\fxfliquid{(4.75,5.75) -- (5,5.75) -- (5,6) -- (4.75,6) -- cycle}
\fxfcell{5}{5}{5.25}{5.25}{1}
\fxfliquid{(5,5) -- (5.25,5) -- (5.25,5.25) -- (5,5.25) -- cycle}
\fxfcell{5.25}{5}{5.5}{5.25}{1}
\fxfliquid{(5.25,5) -- (5.5,5) -- (5.5,5.25) -- (5.25,5.25) -- cycle}
\fxfcell{5}{5.25}{5.25}{5.5}{1}
\fxfliquid{(5,5.25) -- (5.25,5.25) -- (5.25,5.5) -- (5,5.5) -- cycle}
\fxfcell{5.25}{5.25}{5.5}{5.5}{1}
\fxfliquid{(5.25,5.25) -- (5.5,5.25) -- (5.5,5.5) -- (5.25,5.5) -- cycle}
\fxfcell{5}{5.5}{5.25}{5.75}{1}
\fxfliquid{(5,5.5) -- (5.25,5.5) -- (5.25,5.75) -- (5,5.75) -- cycle}
\fxfcell{5.25}{5.5}{5.5}{5.75}{1}
\fxfliquid{(5.25,5.5) -- (5.5,5.5) -- (5.5,5.75) -- (5.25,5.75) -- cycle}
\fxfcell{5}{5.75}{5.25}{6}{1}
\fxfliquid{(5,5.75) -- (5.25,5.75) -- (5.25,6) -- (5,6) -- cycle}
\fxfcell{5.25}{5.75}{5.5}{6}{1}
\fxfliquid{(5.25,5.75) -- (5.5,5.75) -- (5.5,5.980812) -- (5.465599,6) -- (5.25,6) -- cycle}
\fxfplic{5.5}{5.980812}{5.465599}{6}
\fxfcell{4}{6}{4.25}{6.25}{1}
\fxfliquid{(4,6) -- (4.25,6) -- (4.25,6.25) -- (4,6.25) -- cycle}
\fxfcell{4.25}{6}{4.5}{6.25}{1}
\fxfliquid{(4.25,6) -- (4.5,6) -- (4.5,6.25) -- (4.25,6.25) -- cycle}
\fxfcell{4.5}{6}{4.75}{6.25}{1}
\fxfliquid{(4.5,6) -- (4.75,6) -- (4.75,6.25) -- (4.5,6.25) -- cycle}
\fxfcell{4.75}{6}{5}{6.25}{1}
\fxfliquid{(4.75,6) -- (5,6) -- (5,6.219618) -- (4.887192,6.25) -- (4.75,6.25) -- cycle}
\fxfplic{5}{6.219618}{4.887192}{6.25}
\fxfcell{4}{6.25}{4.25}{6.5}{1}
\fxfliquid{(4,6.25) -- (4.25,6.25) -- (4.25,6.331586) -- (4,6.312301) -- cycle}
\fxfplic{4.25}{6.331586}{4}{6.312301}
\fxfcell{4.25}{6.25}{4.5}{6.5}{1}
\fxfliquid{(4.25,6.25) -- (4.5,6.25) -- (4.5,6.323328) -- (4.25,6.331808) -- cycle}
\fxfplic{4.5}{6.323328}{4.25}{6.331808}
\fxfcell{4.5}{6.25}{4.75}{6.5}{1}
\fxfliquid{(4.5,6.25) -- (4.75,6.25) -- (4.75,6.286416) -- (4.5,6.323551) -- cycle}
\fxfplic{4.75}{6.286416}{4.5}{6.323551}
\fxfcell{4.75}{6.25}{5}{6.5}{1}
\fxfliquid{(4.75,6.25) -- (4.887292,6.25) -- (4.75,6.286976) -- cycle}
\fxfplic{4.887292}{6.25}{4.75}{6.286976}
\fxfcell{5}{6}{5.25}{6.25}{1}
\fxfliquid{(5,6) -- (5.25,6) -- (5.25,6.12001) -- (5,6.220636) -- cycle}
\fxfplic{5.25}{6.12001}{5}{6.220636}
\fxfcell{5.25}{6}{5.5}{6.25}{1}
\fxfliquid{(5.25,6) -- (5.468716,6) -- (5.25,6.121998) -- cycle}
\fxfplic{5.468716}{6}{5.25}{6.121998}
\fxfcell{5}{6.25}{5.25}{6.5}{1}
\fxfcell{5.25}{6.25}{5.5}{6.5}{1}
\fxfjump{1}{1}{2}{4.25}{0.5}{0.134521}{0.397586}{0}
\fxfjump{1}{1}{2}{4.75}{0.5}{0.105154}{0.398146}{0}
\fxfjump{1}{1}{2}{5.25}{0.5}{0.035039}{0.398737}{0}
\fxfjump{1}{1}{2.5}{4.25}{0.5}{0.025533}{0.034506}{0}
\fxfjump{1}{1}{2.5}{4.75}{0.5}{0.055233}{0.082967}{0}
\fxfjump{1}{1}{2.5}{5.25}{0.5}{0.124868}{0.257334}{0}
\fxfjump{1}{1}{2.5}{5.75}{0.5}{0.060092}{0.380372}{0}
\fxfjump{1}{1}{3}{5.25}{0.5}{0.001707}{0.002135}{0}
\fxfjump{1}{1}{3}{5.75}{0.5}{0.090339}{0.16412}{0}
\fxfjump{1}{1}{3}{6.25}{0.5}{0.023173}{0.395913}{0}
\fxfjump{1}{1}{3.5}{5.75}{0.5}{0.003669}{0.004616}{0}
\fxfjump{1}{1}{3.5}{6.25}{0.5}{0.059127}{0.223858}{0}
\fxfjump{1}{1}{4}{5.875}{0.25}{0}{0.022437}{1}
\fxfjump{1}{1}{4}{6.125}{0.25}{0.039114}{0.101579}{0}
\fxfjump{1}{1}{4}{6.375}{0.25}{-0.017527}{-0.006382}{0}
\fxfjump{1}{1}{4.25}{6.375}{0.25}{0.001741}{0.028892}{0}
\fxfjump{1}{1}{4.5}{6.375}{0.25}{-0.007123}{-0.133311}{0}
\fxfjump{1}{1}{4.75}{6.125}{0.25}{-0.002163}{-0.010909}{0}
\fxfjump{1}{1}{4.75}{6.375}{0.25}{-0.014142}{-0.536472}{0}
\fxfjump{1}{1}{5}{6.125}{0.25}{-0.022985}{-0.138253}{0}
\fxfjump{1}{1}{5}{6.375}{0.25}{-0.003203}{-0.748703}{0}
\fxfjump{1}{1}{5.25}{5.875}{0.25}{-0.000417}{-0.00208}{0}
\fxfjump{1}{1}{5.25}{6.125}{0.25}{-0.036929}{-0.410157}{0}
\fxfjump{1}{1}{5.5}{5.875}{0.25}{-0.033923}{-0.216249}{0}
\fxfjump{1}{1}{5.5}{6.125}{0.25}{-0.016857}{-0.779232}{0}
\fxfjump{1}{2}{2.25}{4}{0.5}{0.007279}{0.011069}{0}
\fxfjump{1}{2}{2.25}{4.5}{0.5}{-0.029624}{-0.049238}{0}
\fxfjump{1}{2}{2.25}{5}{0.5}{-0.070613}{-0.201073}{0}
\fxfjump{1}{2}{2.75}{5}{0.5}{-0.001707}{-0.002135}{0}
\fxfjump{1}{2}{2.25}{5.5}{0.5}{-0.035039}{-0.398737}{0}
\fxfjump{1}{2}{2.75}{5.5}{0.5}{-0.092031}{-0.165718}{0}
\fxfjump{1}{2}{3.25}{5.5}{0.5}{-0.003669}{-0.004616}{0}
\fxfjump{1}{2}{5.375}{5.75}{0.25}{-0.000417}{-0.00208}{0}
\fxfjump{1}{2}{2.75}{6}{0.5}{-0.060092}{-0.380372}{0}
\fxfjump{1}{2}{3.25}{6}{0.5}{-0.134374}{-0.297947}{0}
\fxfjump{1}{2}{3.75}{6}{0.5}{-0.078227}{-0.119258}{0}
\fxfjump{1}{2}{4.125}{6}{0.25}{0}{-0.015671}{1}
\fxfjump{1}{2}{4.875}{6}{0.25}{-0.002163}{-0.010909}{0}
\fxfjump{1}{2}{5.125}{6}{0.25}{-0.02515}{-0.14881}{0}
\fxfjump{1}{2}{5.375}{6}{0.25}{-0.061696}{-0.507764}{0}
\fxfjump{1}{2}{4.125}{6.25}{0.25}{-0.054001}{-0.405269}{0}
\fxfjump{1}{2}{4.375}{6.25}{0.25}{-0.054385}{-0.412772}{0}
\fxfjump{1}{2}{4.625}{6.25}{0.25}{-0.061514}{-0.501399}{0}
\fxfjump{1}{2}{4.875}{6.25}{0.25}{-0.073511}{-0.721207}{0}
\fxfjump{1}{2}{5.125}{6.25}{0.25}{-0.053762}{-0.783653}{0}
\fxfjump{1}{2}{5.375}{6.25}{0.25}{-0.016857}{-0.779232}{0}
\fxfjump{1}{2}{3.25}{6.5}{0.5}{-0.023173}{-0.395913}{0}
\fxfjump{1}{2}{3.75}{6.5}{0.5}{-0.082051}{-0.408191}{0}
\fxfjump{1}{2}{4.125}{6.5}{0.25}{-0.021819}{-0.718584}{0}
\fxfjump{1}{2}{4.375}{6.5}{0.25}{-0.024465}{-0.780304}{0}
\fxfjump{1}{2}{4.625}{6.5}{0.25}{-0.017343}{-0.778333}{0}
\fxfjump{1}{2}{4.875}{6.5}{0.25}{-0.003203}{-0.748703}{0}
\fxfjump{2}{1}{2}{4.25}{0.5}{0.160054}{0.214311}{0}
\fxfjump{2}{1}{2}{4.75}{0.5}{0.160387}{0.27512}{0}
\fxfjump{2}{1}{2.5}{5.25}{0.5}{0.161615}{0.15089}{0}
\fxfjump{2}{1}{3}{5.75}{0.5}{0.154197}{0.12691}{0}
\fxfjump{2}{1}{3.5}{6.25}{0.5}{0.160588}{0.275447}{0}
\fxfjump{2}{1}{4}{6.125}{0.25}{0}{-0.035963}{1}
\fxfjump{2}{1}{4}{6.375}{0.25}{-0.078193}{-0.175936}{0}
\fxfjump{2}{1}{5.25}{6.125}{0.25}{-0.078937}{-0.397193}{0}
\fxfjump{2}{2}{2.25}{5}{0.5}{-0.160914}{-0.207587}{0}
\fxfjump{2}{2}{2.75}{5.5}{0.5}{-0.153881}{-0.125533}{0}
\fxfjump{2}{2}{3.25}{6}{0.5}{-0.161216}{-0.161947}{0}
\fxfjump{2}{2}{5.375}{6}{0.25}{-0.07897}{-0.294446}{0}
\fxfjump{2}{2}{4.125}{6.25}{0.25}{-0.075821}{-0.246759}{0}
\fxfjump{2}{2}{4.375}{6.25}{0.25}{-0.078849}{-0.271125}{0}
\fxfjump{2}{2}{4.625}{6.25}{0.25}{-0.078857}{-0.291199}{0}
\fxfjump{2}{2}{4.875}{6.25}{0.25}{-0.078877}{-0.350587}{0}
\fxfjump{2}{2}{5.125}{6.25}{0.25}{-0.078911}{-0.521108}{0}
\fxfjump{2}{2}{3.75}{6.5}{0.5}{-0.160278}{-0.33277}{0}
\fxfjump{2}{2}{4.125}{6.5}{0.25}{0}{-0.067844}{1}
\fxfhanging{4}{4.25}
\fxfhanging{4}{4.75}
\fxfhanging{4}{5.25}
\fxfhanging{4}{5.75}
\fxfhanging{4}{6.25}
%
      \fxfnone
      \def\fxfcell##1##2##3##4##5{\draw[tr cell] (##1,##2) rectangle (##3,##4);}%
      \def\fxfplic##1##2##3##4{\draw[tr plic] (##1,##2) -- (##3,##4);}%
      % Generated by docs/techref/figures/generate.jl; do not edit.
\fxfcell{2}{4}{2.5}{4.5}{0}
\fxfliquid{(2.062341,4) -- (2.5,4) -- (2.5,4.5) -- (2.097185,4.5) -- cycle}
\fxfplic{2.062341}{4}{2.097185}{4.5}
\fxfcell{2.5}{4}{3}{4.5}{0}
\fxfliquid{(2.5,4) -- (3,4) -- (3,4.5) -- (2.5,4.5) -- cycle}
\fxfcell{3}{4}{3.5}{4.5}{0}
\fxfliquid{(3,4) -- (3.5,4) -- (3.5,4.5) -- (3,4.5) -- cycle}
\fxfcell{3.5}{4}{4}{4.5}{0}
\fxfliquid{(3.5,4) -- (4,4) -- (4,4.5) -- (3.5,4.5) -- cycle}
\fxfcell{2}{4.5}{2.5}{5}{0}
\fxfliquid{(2.092907,4.5) -- (2.5,4.5) -- (2.5,5) -- (2.251465,5) -- cycle}
\fxfplic{2.092907}{4.5}{2.251465}{5}
\fxfcell{2.5}{4.5}{3}{5}{0}
\fxfliquid{(2.5,4.5) -- (3,4.5) -- (3,5) -- (2.5,5) -- cycle}
\fxfcell{3}{4.5}{3.5}{5}{0}
\fxfliquid{(3,4.5) -- (3.5,4.5) -- (3.5,5) -- (3,5) -- cycle}
\fxfcell{3.5}{4.5}{4}{5}{0}
\fxfliquid{(3.5,4.5) -- (4,4.5) -- (4,5) -- (3.5,5) -- cycle}
\fxfcell{2}{5}{2.5}{5.5}{0}
\fxfliquid{(2.235883,5) -- (2.5,5) -- (2.5,5.410432) -- cycle}
\fxfplic{2.235883}{5}{2.5}{5.410432}
\fxfcell{2.5}{5}{3}{5.5}{0}
\fxfliquid{(2.5,5) -- (3,5) -- (3,5.5) -- (2.558302,5.5) -- (2.5,5.4094) -- cycle}
\fxfplic{2.558302}{5.5}{2.5}{5.4094}
\fxfcell{3}{5}{3.5}{5.5}{0}
\fxfliquid{(3,5) -- (3.5,5) -- (3.5,5.5) -- (3,5.5) -- cycle}
\fxfcell{3.5}{5}{4}{5.5}{0}
\fxfliquid{(3.5,5) -- (4,5) -- (4,5.5) -- (3.5,5.5) -- cycle}
\fxfcell{2}{5.5}{2.5}{6}{0}
\fxfcell{2.5}{5.5}{3}{6}{0}
\fxfliquid{(2.5,5.5) -- (3,5.5) -- (3,5.880728) -- (2.5,5.510641) -- cycle}
\fxfplic{3}{5.880728}{2.5}{5.510641}
\fxfcell{3}{5.5}{3.5}{6}{0}
\fxfliquid{(3,5.5) -- (3.5,5.5) -- (3.5,6) -- (3.141854,6) -- (3,5.919773) -- cycle}
\fxfplic{3.141854}{6}{3}{5.919773}
\fxfcell{3.5}{5.5}{4}{6}{0}
\fxfliquid{(3.5,5.5) -- (4,5.5) -- (4,6) -- (3.5,6) -- cycle}
\fxfcell{2}{6}{2.5}{6.5}{0}
\fxfcell{2.5}{6}{3}{6.5}{0}
\fxfcell{3}{6}{3.5}{6.5}{0}
\fxfliquid{(3.143521,6) -- (3.5,6) -- (3.5,6.201611) -- cycle}
\fxfplic{3.143521}{6}{3.5}{6.201611}
\fxfcell{3.5}{6}{4}{6.5}{0}
\fxfliquid{(3.5,6) -- (4,6) -- (4,6.322031) -- (3.5,6.189898) -- cycle}
\fxfplic{4}{6.322031}{3.5}{6.189898}
\fxfcell{4}{4}{4.25}{4.25}{1}
\fxfliquid{(4,4) -- (4.25,4) -- (4.25,4.25) -- (4,4.25) -- cycle}
\fxfcell{4.25}{4}{4.5}{4.25}{1}
\fxfliquid{(4.25,4) -- (4.5,4) -- (4.5,4.25) -- (4.25,4.25) -- cycle}
\fxfcell{4.5}{4}{4.75}{4.25}{1}
\fxfliquid{(4.5,4) -- (4.75,4) -- (4.75,4.25) -- (4.5,4.25) -- cycle}
\fxfcell{4.75}{4}{5}{4.25}{1}
\fxfliquid{(4.75,4) -- (5,4) -- (5,4.25) -- (4.75,4.25) -- cycle}
\fxfcell{4}{4.25}{4.25}{4.5}{1}
\fxfliquid{(4,4.25) -- (4.25,4.25) -- (4.25,4.5) -- (4,4.5) -- cycle}
\fxfcell{4.25}{4.25}{4.5}{4.5}{1}
\fxfliquid{(4.25,4.25) -- (4.5,4.25) -- (4.5,4.5) -- (4.25,4.5) -- cycle}
\fxfcell{4.5}{4.25}{4.75}{4.5}{1}
\fxfliquid{(4.5,4.25) -- (4.75,4.25) -- (4.75,4.5) -- (4.5,4.5) -- cycle}
\fxfcell{4.75}{4.25}{5}{4.5}{1}
\fxfliquid{(4.75,4.25) -- (5,4.25) -- (5,4.5) -- (4.75,4.5) -- cycle}
\fxfcell{4}{4.5}{4.25}{4.75}{1}
\fxfliquid{(4,4.5) -- (4.25,4.5) -- (4.25,4.75) -- (4,4.75) -- cycle}
\fxfcell{4.25}{4.5}{4.5}{4.75}{1}
\fxfliquid{(4.25,4.5) -- (4.5,4.5) -- (4.5,4.75) -- (4.25,4.75) -- cycle}
\fxfcell{4.5}{4.5}{4.75}{4.75}{1}
\fxfliquid{(4.5,4.5) -- (4.75,4.5) -- (4.75,4.75) -- (4.5,4.75) -- cycle}
\fxfcell{4.75}{4.5}{5}{4.75}{1}
\fxfliquid{(4.75,4.5) -- (5,4.5) -- (5,4.75) -- (4.75,4.75) -- cycle}
\fxfcell{4}{4.75}{4.25}{5}{1}
\fxfliquid{(4,4.75) -- (4.25,4.75) -- (4.25,5) -- (4,5) -- cycle}
\fxfcell{4.25}{4.75}{4.5}{5}{1}
\fxfliquid{(4.25,4.75) -- (4.5,4.75) -- (4.5,5) -- (4.25,5) -- cycle}
\fxfcell{4.5}{4.75}{4.75}{5}{1}
\fxfliquid{(4.5,4.75) -- (4.75,4.75) -- (4.75,5) -- (4.5,5) -- cycle}
\fxfcell{4.75}{4.75}{5}{5}{1}
\fxfliquid{(4.75,4.75) -- (5,4.75) -- (5,5) -- (4.75,5) -- cycle}
\fxfcell{5}{4}{5.25}{4.25}{1}
\fxfliquid{(5,4) -- (5.25,4) -- (5.25,4.25) -- (5,4.25) -- cycle}
\fxfcell{5.25}{4}{5.5}{4.25}{1}
\fxfliquid{(5.25,4) -- (5.5,4) -- (5.5,4.25) -- (5.25,4.25) -- cycle}
\fxfcell{5}{4.25}{5.25}{4.5}{1}
\fxfliquid{(5,4.25) -- (5.25,4.25) -- (5.25,4.5) -- (5,4.5) -- cycle}
\fxfcell{5.25}{4.25}{5.5}{4.5}{1}
\fxfliquid{(5.25,4.25) -- (5.5,4.25) -- (5.5,4.5) -- (5.25,4.5) -- cycle}
\fxfcell{5}{4.5}{5.25}{4.75}{1}
\fxfliquid{(5,4.5) -- (5.25,4.5) -- (5.25,4.75) -- (5,4.75) -- cycle}
\fxfcell{5.25}{4.5}{5.5}{4.75}{1}
\fxfliquid{(5.25,4.5) -- (5.5,4.5) -- (5.5,4.75) -- (5.25,4.75) -- cycle}
\fxfcell{5}{4.75}{5.25}{5}{1}
\fxfliquid{(5,4.75) -- (5.25,4.75) -- (5.25,5) -- (5,5) -- cycle}
\fxfcell{5.25}{4.75}{5.5}{5}{1}
\fxfliquid{(5.25,4.75) -- (5.5,4.75) -- (5.5,5) -- (5.25,5) -- cycle}
\fxfcell{4}{5}{4.25}{5.25}{1}
\fxfliquid{(4,5) -- (4.25,5) -- (4.25,5.25) -- (4,5.25) -- cycle}
\fxfcell{4.25}{5}{4.5}{5.25}{1}
\fxfliquid{(4.25,5) -- (4.5,5) -- (4.5,5.25) -- (4.25,5.25) -- cycle}
\fxfcell{4.5}{5}{4.75}{5.25}{1}
\fxfliquid{(4.5,5) -- (4.75,5) -- (4.75,5.25) -- (4.5,5.25) -- cycle}
\fxfcell{4.75}{5}{5}{5.25}{1}
\fxfliquid{(4.75,5) -- (5,5) -- (5,5.25) -- (4.75,5.25) -- cycle}
\fxfcell{4}{5.25}{4.25}{5.5}{1}
\fxfliquid{(4,5.25) -- (4.25,5.25) -- (4.25,5.5) -- (4,5.5) -- cycle}
\fxfcell{4.25}{5.25}{4.5}{5.5}{1}
\fxfliquid{(4.25,5.25) -- (4.5,5.25) -- (4.5,5.5) -- (4.25,5.5) -- cycle}
\fxfcell{4.5}{5.25}{4.75}{5.5}{1}
\fxfliquid{(4.5,5.25) -- (4.75,5.25) -- (4.75,5.5) -- (4.5,5.5) -- cycle}
\fxfcell{4.75}{5.25}{5}{5.5}{1}
\fxfliquid{(4.75,5.25) -- (5,5.25) -- (5,5.5) -- (4.75,5.5) -- cycle}
\fxfcell{4}{5.5}{4.25}{5.75}{1}
\fxfliquid{(4,5.5) -- (4.25,5.5) -- (4.25,5.75) -- (4,5.75) -- cycle}
\fxfcell{4.25}{5.5}{4.5}{5.75}{1}
\fxfliquid{(4.25,5.5) -- (4.5,5.5) -- (4.5,5.75) -- (4.25,5.75) -- cycle}
\fxfcell{4.5}{5.5}{4.75}{5.75}{1}
\fxfliquid{(4.5,5.5) -- (4.75,5.5) -- (4.75,5.75) -- (4.5,5.75) -- cycle}
\fxfcell{4.75}{5.5}{5}{5.75}{1}
\fxfliquid{(4.75,5.5) -- (5,5.5) -- (5,5.75) -- (4.75,5.75) -- cycle}
\fxfcell{4}{5.75}{4.25}{6}{1}
\fxfliquid{(4,5.75) -- (4.25,5.75) -- (4.25,6) -- (4,6) -- cycle}
\fxfcell{4.25}{5.75}{4.5}{6}{1}
\fxfliquid{(4.25,5.75) -- (4.5,5.75) -- (4.5,6) -- (4.25,6) -- cycle}
\fxfcell{4.5}{5.75}{4.75}{6}{1}
\fxfliquid{(4.5,5.75) -- (4.75,5.75) -- (4.75,6) -- (4.5,6) -- cycle}
\fxfcell{4.75}{5.75}{5}{6}{1}
\fxfliquid{(4.75,5.75) -- (5,5.75) -- (5,6) -- (4.75,6) -- cycle}
\fxfcell{5}{5}{5.25}{5.25}{1}
\fxfliquid{(5,5) -- (5.25,5) -- (5.25,5.25) -- (5,5.25) -- cycle}
\fxfcell{5.25}{5}{5.5}{5.25}{1}
\fxfliquid{(5.25,5) -- (5.5,5) -- (5.5,5.25) -- (5.25,5.25) -- cycle}
\fxfcell{5}{5.25}{5.25}{5.5}{1}
\fxfliquid{(5,5.25) -- (5.25,5.25) -- (5.25,5.5) -- (5,5.5) -- cycle}
\fxfcell{5.25}{5.25}{5.5}{5.5}{1}
\fxfliquid{(5.25,5.25) -- (5.5,5.25) -- (5.5,5.5) -- (5.25,5.5) -- cycle}
\fxfcell{5}{5.5}{5.25}{5.75}{1}
\fxfliquid{(5,5.5) -- (5.25,5.5) -- (5.25,5.75) -- (5,5.75) -- cycle}
\fxfcell{5.25}{5.5}{5.5}{5.75}{1}
\fxfliquid{(5.25,5.5) -- (5.5,5.5) -- (5.5,5.75) -- (5.25,5.75) -- cycle}
\fxfcell{5}{5.75}{5.25}{6}{1}
\fxfliquid{(5,5.75) -- (5.25,5.75) -- (5.25,6) -- (5,6) -- cycle}
\fxfcell{5.25}{5.75}{5.5}{6}{1}
\fxfliquid{(5.25,5.75) -- (5.5,5.75) -- (5.5,5.980812) -- (5.465599,6) -- (5.25,6) -- cycle}
\fxfplic{5.5}{5.980812}{5.465599}{6}
\fxfcell{4}{6}{4.25}{6.25}{1}
\fxfliquid{(4,6) -- (4.25,6) -- (4.25,6.25) -- (4,6.25) -- cycle}
\fxfcell{4.25}{6}{4.5}{6.25}{1}
\fxfliquid{(4.25,6) -- (4.5,6) -- (4.5,6.25) -- (4.25,6.25) -- cycle}
\fxfcell{4.5}{6}{4.75}{6.25}{1}
\fxfliquid{(4.5,6) -- (4.75,6) -- (4.75,6.25) -- (4.5,6.25) -- cycle}
\fxfcell{4.75}{6}{5}{6.25}{1}
\fxfliquid{(4.75,6) -- (5,6) -- (5,6.219618) -- (4.887192,6.25) -- (4.75,6.25) -- cycle}
\fxfplic{5}{6.219618}{4.887192}{6.25}
\fxfcell{4}{6.25}{4.25}{6.5}{1}
\fxfliquid{(4,6.25) -- (4.25,6.25) -- (4.25,6.331586) -- (4,6.312301) -- cycle}
\fxfplic{4.25}{6.331586}{4}{6.312301}
\fxfcell{4.25}{6.25}{4.5}{6.5}{1}
\fxfliquid{(4.25,6.25) -- (4.5,6.25) -- (4.5,6.323328) -- (4.25,6.331808) -- cycle}
\fxfplic{4.5}{6.323328}{4.25}{6.331808}
\fxfcell{4.5}{6.25}{4.75}{6.5}{1}
\fxfliquid{(4.5,6.25) -- (4.75,6.25) -- (4.75,6.286416) -- (4.5,6.323551) -- cycle}
\fxfplic{4.75}{6.286416}{4.5}{6.323551}
\fxfcell{4.75}{6.25}{5}{6.5}{1}
\fxfliquid{(4.75,6.25) -- (4.887292,6.25) -- (4.75,6.286976) -- cycle}
\fxfplic{4.887292}{6.25}{4.75}{6.286976}
\fxfcell{5}{6}{5.25}{6.25}{1}
\fxfliquid{(5,6) -- (5.25,6) -- (5.25,6.12001) -- (5,6.220636) -- cycle}
\fxfplic{5.25}{6.12001}{5}{6.220636}
\fxfcell{5.25}{6}{5.5}{6.25}{1}
\fxfliquid{(5.25,6) -- (5.468716,6) -- (5.25,6.121998) -- cycle}
\fxfplic{5.468716}{6}{5.25}{6.121998}
\fxfcell{5}{6.25}{5.25}{6.5}{1}
\fxfcell{5.25}{6.25}{5.5}{6.5}{1}
\fxfjump{1}{1}{2}{4.25}{0.5}{0.134521}{0.397586}{0}
\fxfjump{1}{1}{2}{4.75}{0.5}{0.105154}{0.398146}{0}
\fxfjump{1}{1}{2}{5.25}{0.5}{0.035039}{0.398737}{0}
\fxfjump{1}{1}{2.5}{4.25}{0.5}{0.025533}{0.034506}{0}
\fxfjump{1}{1}{2.5}{4.75}{0.5}{0.055233}{0.082967}{0}
\fxfjump{1}{1}{2.5}{5.25}{0.5}{0.124868}{0.257334}{0}
\fxfjump{1}{1}{2.5}{5.75}{0.5}{0.060092}{0.380372}{0}
\fxfjump{1}{1}{3}{5.25}{0.5}{0.001707}{0.002135}{0}
\fxfjump{1}{1}{3}{5.75}{0.5}{0.090339}{0.16412}{0}
\fxfjump{1}{1}{3}{6.25}{0.5}{0.023173}{0.395913}{0}
\fxfjump{1}{1}{3.5}{5.75}{0.5}{0.003669}{0.004616}{0}
\fxfjump{1}{1}{3.5}{6.25}{0.5}{0.059127}{0.223858}{0}
\fxfjump{1}{1}{4}{5.875}{0.25}{0}{0.022437}{1}
\fxfjump{1}{1}{4}{6.125}{0.25}{0.039114}{0.101579}{0}
\fxfjump{1}{1}{4}{6.375}{0.25}{-0.017527}{-0.006382}{0}
\fxfjump{1}{1}{4.25}{6.375}{0.25}{0.001741}{0.028892}{0}
\fxfjump{1}{1}{4.5}{6.375}{0.25}{-0.007123}{-0.133311}{0}
\fxfjump{1}{1}{4.75}{6.125}{0.25}{-0.002163}{-0.010909}{0}
\fxfjump{1}{1}{4.75}{6.375}{0.25}{-0.014142}{-0.536472}{0}
\fxfjump{1}{1}{5}{6.125}{0.25}{-0.022985}{-0.138253}{0}
\fxfjump{1}{1}{5}{6.375}{0.25}{-0.003203}{-0.748703}{0}
\fxfjump{1}{1}{5.25}{5.875}{0.25}{-0.000417}{-0.00208}{0}
\fxfjump{1}{1}{5.25}{6.125}{0.25}{-0.036929}{-0.410157}{0}
\fxfjump{1}{1}{5.5}{5.875}{0.25}{-0.033923}{-0.216249}{0}
\fxfjump{1}{1}{5.5}{6.125}{0.25}{-0.016857}{-0.779232}{0}
\fxfjump{1}{2}{2.25}{4}{0.5}{0.007279}{0.011069}{0}
\fxfjump{1}{2}{2.25}{4.5}{0.5}{-0.029624}{-0.049238}{0}
\fxfjump{1}{2}{2.25}{5}{0.5}{-0.070613}{-0.201073}{0}
\fxfjump{1}{2}{2.75}{5}{0.5}{-0.001707}{-0.002135}{0}
\fxfjump{1}{2}{2.25}{5.5}{0.5}{-0.035039}{-0.398737}{0}
\fxfjump{1}{2}{2.75}{5.5}{0.5}{-0.092031}{-0.165718}{0}
\fxfjump{1}{2}{3.25}{5.5}{0.5}{-0.003669}{-0.004616}{0}
\fxfjump{1}{2}{5.375}{5.75}{0.25}{-0.000417}{-0.00208}{0}
\fxfjump{1}{2}{2.75}{6}{0.5}{-0.060092}{-0.380372}{0}
\fxfjump{1}{2}{3.25}{6}{0.5}{-0.134374}{-0.297947}{0}
\fxfjump{1}{2}{3.75}{6}{0.5}{-0.078227}{-0.119258}{0}
\fxfjump{1}{2}{4.125}{6}{0.25}{0}{-0.015671}{1}
\fxfjump{1}{2}{4.875}{6}{0.25}{-0.002163}{-0.010909}{0}
\fxfjump{1}{2}{5.125}{6}{0.25}{-0.02515}{-0.14881}{0}
\fxfjump{1}{2}{5.375}{6}{0.25}{-0.061696}{-0.507764}{0}
\fxfjump{1}{2}{4.125}{6.25}{0.25}{-0.054001}{-0.405269}{0}
\fxfjump{1}{2}{4.375}{6.25}{0.25}{-0.054385}{-0.412772}{0}
\fxfjump{1}{2}{4.625}{6.25}{0.25}{-0.061514}{-0.501399}{0}
\fxfjump{1}{2}{4.875}{6.25}{0.25}{-0.073511}{-0.721207}{0}
\fxfjump{1}{2}{5.125}{6.25}{0.25}{-0.053762}{-0.783653}{0}
\fxfjump{1}{2}{5.375}{6.25}{0.25}{-0.016857}{-0.779232}{0}
\fxfjump{1}{2}{3.25}{6.5}{0.5}{-0.023173}{-0.395913}{0}
\fxfjump{1}{2}{3.75}{6.5}{0.5}{-0.082051}{-0.408191}{0}
\fxfjump{1}{2}{4.125}{6.5}{0.25}{-0.021819}{-0.718584}{0}
\fxfjump{1}{2}{4.375}{6.5}{0.25}{-0.024465}{-0.780304}{0}
\fxfjump{1}{2}{4.625}{6.5}{0.25}{-0.017343}{-0.778333}{0}
\fxfjump{1}{2}{4.875}{6.5}{0.25}{-0.003203}{-0.748703}{0}
\fxfjump{2}{1}{2}{4.25}{0.5}{0.160054}{0.214311}{0}
\fxfjump{2}{1}{2}{4.75}{0.5}{0.160387}{0.27512}{0}
\fxfjump{2}{1}{2.5}{5.25}{0.5}{0.161615}{0.15089}{0}
\fxfjump{2}{1}{3}{5.75}{0.5}{0.154197}{0.12691}{0}
\fxfjump{2}{1}{3.5}{6.25}{0.5}{0.160588}{0.275447}{0}
\fxfjump{2}{1}{4}{6.125}{0.25}{0}{-0.035963}{1}
\fxfjump{2}{1}{4}{6.375}{0.25}{-0.078193}{-0.175936}{0}
\fxfjump{2}{1}{5.25}{6.125}{0.25}{-0.078937}{-0.397193}{0}
\fxfjump{2}{2}{2.25}{5}{0.5}{-0.160914}{-0.207587}{0}
\fxfjump{2}{2}{2.75}{5.5}{0.5}{-0.153881}{-0.125533}{0}
\fxfjump{2}{2}{3.25}{6}{0.5}{-0.161216}{-0.161947}{0}
\fxfjump{2}{2}{5.375}{6}{0.25}{-0.07897}{-0.294446}{0}
\fxfjump{2}{2}{4.125}{6.25}{0.25}{-0.075821}{-0.246759}{0}
\fxfjump{2}{2}{4.375}{6.25}{0.25}{-0.078849}{-0.271125}{0}
\fxfjump{2}{2}{4.625}{6.25}{0.25}{-0.078857}{-0.291199}{0}
\fxfjump{2}{2}{4.875}{6.25}{0.25}{-0.078877}{-0.350587}{0}
\fxfjump{2}{2}{5.125}{6.25}{0.25}{-0.078911}{-0.521108}{0}
\fxfjump{2}{2}{3.75}{6.5}{0.5}{-0.160278}{-0.33277}{0}
\fxfjump{2}{2}{4.125}{6.5}{0.25}{0}{-0.067844}{1}
\fxfhanging{4}{4.25}
\fxfhanging{4}{4.75}
\fxfhanging{4}{5.25}
\fxfhanging{4}{5.75}
\fxfhanging{4}{6.25}
%
      \fxfnone
      % Faces with a jump: thick red; accelerations: arrows along the face
      % normal, length 0.45 per unit of the largest acceleration (black with
      % a jump, purple without: coupling through the metric).
      \def\fxfjump##1##2##3##4##5##6##7##8{\ifnum##1=#1
        \ifnum##8=0
          \ifnum##2=1 \draw[line width=1.6pt,trred] (##3,##4-##5/2+0.03) -- (##3,##4+##5/2-0.03);
          \else \draw[line width=1.6pt,trred] (##3-##5/2+0.03,##4) -- (##3+##5/2-0.03,##4);\fi
          \def\fxfarrow{black}%
        \else \def\fxfarrow{trpurple}\fi
        \ifnum##2=1 \draw[tr thin vector,\fxfarrow,line width=0.7pt] (##3,##4) -- ++(0.45*##7,0);
        \else \draw[tr thin vector,\fxfarrow,line width=0.7pt] (##3,##4) -- ++(0,0.45*##7);\fi\fi}%
      % Generated by docs/techref/figures/generate.jl; do not edit.
\fxfcell{2}{4}{2.5}{4.5}{0}
\fxfliquid{(2.062341,4) -- (2.5,4) -- (2.5,4.5) -- (2.097185,4.5) -- cycle}
\fxfplic{2.062341}{4}{2.097185}{4.5}
\fxfcell{2.5}{4}{3}{4.5}{0}
\fxfliquid{(2.5,4) -- (3,4) -- (3,4.5) -- (2.5,4.5) -- cycle}
\fxfcell{3}{4}{3.5}{4.5}{0}
\fxfliquid{(3,4) -- (3.5,4) -- (3.5,4.5) -- (3,4.5) -- cycle}
\fxfcell{3.5}{4}{4}{4.5}{0}
\fxfliquid{(3.5,4) -- (4,4) -- (4,4.5) -- (3.5,4.5) -- cycle}
\fxfcell{2}{4.5}{2.5}{5}{0}
\fxfliquid{(2.092907,4.5) -- (2.5,4.5) -- (2.5,5) -- (2.251465,5) -- cycle}
\fxfplic{2.092907}{4.5}{2.251465}{5}
\fxfcell{2.5}{4.5}{3}{5}{0}
\fxfliquid{(2.5,4.5) -- (3,4.5) -- (3,5) -- (2.5,5) -- cycle}
\fxfcell{3}{4.5}{3.5}{5}{0}
\fxfliquid{(3,4.5) -- (3.5,4.5) -- (3.5,5) -- (3,5) -- cycle}
\fxfcell{3.5}{4.5}{4}{5}{0}
\fxfliquid{(3.5,4.5) -- (4,4.5) -- (4,5) -- (3.5,5) -- cycle}
\fxfcell{2}{5}{2.5}{5.5}{0}
\fxfliquid{(2.235883,5) -- (2.5,5) -- (2.5,5.410432) -- cycle}
\fxfplic{2.235883}{5}{2.5}{5.410432}
\fxfcell{2.5}{5}{3}{5.5}{0}
\fxfliquid{(2.5,5) -- (3,5) -- (3,5.5) -- (2.558302,5.5) -- (2.5,5.4094) -- cycle}
\fxfplic{2.558302}{5.5}{2.5}{5.4094}
\fxfcell{3}{5}{3.5}{5.5}{0}
\fxfliquid{(3,5) -- (3.5,5) -- (3.5,5.5) -- (3,5.5) -- cycle}
\fxfcell{3.5}{5}{4}{5.5}{0}
\fxfliquid{(3.5,5) -- (4,5) -- (4,5.5) -- (3.5,5.5) -- cycle}
\fxfcell{2}{5.5}{2.5}{6}{0}
\fxfcell{2.5}{5.5}{3}{6}{0}
\fxfliquid{(2.5,5.5) -- (3,5.5) -- (3,5.880728) -- (2.5,5.510641) -- cycle}
\fxfplic{3}{5.880728}{2.5}{5.510641}
\fxfcell{3}{5.5}{3.5}{6}{0}
\fxfliquid{(3,5.5) -- (3.5,5.5) -- (3.5,6) -- (3.141854,6) -- (3,5.919773) -- cycle}
\fxfplic{3.141854}{6}{3}{5.919773}
\fxfcell{3.5}{5.5}{4}{6}{0}
\fxfliquid{(3.5,5.5) -- (4,5.5) -- (4,6) -- (3.5,6) -- cycle}
\fxfcell{2}{6}{2.5}{6.5}{0}
\fxfcell{2.5}{6}{3}{6.5}{0}
\fxfcell{3}{6}{3.5}{6.5}{0}
\fxfliquid{(3.143521,6) -- (3.5,6) -- (3.5,6.201611) -- cycle}
\fxfplic{3.143521}{6}{3.5}{6.201611}
\fxfcell{3.5}{6}{4}{6.5}{0}
\fxfliquid{(3.5,6) -- (4,6) -- (4,6.322031) -- (3.5,6.189898) -- cycle}
\fxfplic{4}{6.322031}{3.5}{6.189898}
\fxfcell{4}{4}{4.25}{4.25}{1}
\fxfliquid{(4,4) -- (4.25,4) -- (4.25,4.25) -- (4,4.25) -- cycle}
\fxfcell{4.25}{4}{4.5}{4.25}{1}
\fxfliquid{(4.25,4) -- (4.5,4) -- (4.5,4.25) -- (4.25,4.25) -- cycle}
\fxfcell{4.5}{4}{4.75}{4.25}{1}
\fxfliquid{(4.5,4) -- (4.75,4) -- (4.75,4.25) -- (4.5,4.25) -- cycle}
\fxfcell{4.75}{4}{5}{4.25}{1}
\fxfliquid{(4.75,4) -- (5,4) -- (5,4.25) -- (4.75,4.25) -- cycle}
\fxfcell{4}{4.25}{4.25}{4.5}{1}
\fxfliquid{(4,4.25) -- (4.25,4.25) -- (4.25,4.5) -- (4,4.5) -- cycle}
\fxfcell{4.25}{4.25}{4.5}{4.5}{1}
\fxfliquid{(4.25,4.25) -- (4.5,4.25) -- (4.5,4.5) -- (4.25,4.5) -- cycle}
\fxfcell{4.5}{4.25}{4.75}{4.5}{1}
\fxfliquid{(4.5,4.25) -- (4.75,4.25) -- (4.75,4.5) -- (4.5,4.5) -- cycle}
\fxfcell{4.75}{4.25}{5}{4.5}{1}
\fxfliquid{(4.75,4.25) -- (5,4.25) -- (5,4.5) -- (4.75,4.5) -- cycle}
\fxfcell{4}{4.5}{4.25}{4.75}{1}
\fxfliquid{(4,4.5) -- (4.25,4.5) -- (4.25,4.75) -- (4,4.75) -- cycle}
\fxfcell{4.25}{4.5}{4.5}{4.75}{1}
\fxfliquid{(4.25,4.5) -- (4.5,4.5) -- (4.5,4.75) -- (4.25,4.75) -- cycle}
\fxfcell{4.5}{4.5}{4.75}{4.75}{1}
\fxfliquid{(4.5,4.5) -- (4.75,4.5) -- (4.75,4.75) -- (4.5,4.75) -- cycle}
\fxfcell{4.75}{4.5}{5}{4.75}{1}
\fxfliquid{(4.75,4.5) -- (5,4.5) -- (5,4.75) -- (4.75,4.75) -- cycle}
\fxfcell{4}{4.75}{4.25}{5}{1}
\fxfliquid{(4,4.75) -- (4.25,4.75) -- (4.25,5) -- (4,5) -- cycle}
\fxfcell{4.25}{4.75}{4.5}{5}{1}
\fxfliquid{(4.25,4.75) -- (4.5,4.75) -- (4.5,5) -- (4.25,5) -- cycle}
\fxfcell{4.5}{4.75}{4.75}{5}{1}
\fxfliquid{(4.5,4.75) -- (4.75,4.75) -- (4.75,5) -- (4.5,5) -- cycle}
\fxfcell{4.75}{4.75}{5}{5}{1}
\fxfliquid{(4.75,4.75) -- (5,4.75) -- (5,5) -- (4.75,5) -- cycle}
\fxfcell{5}{4}{5.25}{4.25}{1}
\fxfliquid{(5,4) -- (5.25,4) -- (5.25,4.25) -- (5,4.25) -- cycle}
\fxfcell{5.25}{4}{5.5}{4.25}{1}
\fxfliquid{(5.25,4) -- (5.5,4) -- (5.5,4.25) -- (5.25,4.25) -- cycle}
\fxfcell{5}{4.25}{5.25}{4.5}{1}
\fxfliquid{(5,4.25) -- (5.25,4.25) -- (5.25,4.5) -- (5,4.5) -- cycle}
\fxfcell{5.25}{4.25}{5.5}{4.5}{1}
\fxfliquid{(5.25,4.25) -- (5.5,4.25) -- (5.5,4.5) -- (5.25,4.5) -- cycle}
\fxfcell{5}{4.5}{5.25}{4.75}{1}
\fxfliquid{(5,4.5) -- (5.25,4.5) -- (5.25,4.75) -- (5,4.75) -- cycle}
\fxfcell{5.25}{4.5}{5.5}{4.75}{1}
\fxfliquid{(5.25,4.5) -- (5.5,4.5) -- (5.5,4.75) -- (5.25,4.75) -- cycle}
\fxfcell{5}{4.75}{5.25}{5}{1}
\fxfliquid{(5,4.75) -- (5.25,4.75) -- (5.25,5) -- (5,5) -- cycle}
\fxfcell{5.25}{4.75}{5.5}{5}{1}
\fxfliquid{(5.25,4.75) -- (5.5,4.75) -- (5.5,5) -- (5.25,5) -- cycle}
\fxfcell{4}{5}{4.25}{5.25}{1}
\fxfliquid{(4,5) -- (4.25,5) -- (4.25,5.25) -- (4,5.25) -- cycle}
\fxfcell{4.25}{5}{4.5}{5.25}{1}
\fxfliquid{(4.25,5) -- (4.5,5) -- (4.5,5.25) -- (4.25,5.25) -- cycle}
\fxfcell{4.5}{5}{4.75}{5.25}{1}
\fxfliquid{(4.5,5) -- (4.75,5) -- (4.75,5.25) -- (4.5,5.25) -- cycle}
\fxfcell{4.75}{5}{5}{5.25}{1}
\fxfliquid{(4.75,5) -- (5,5) -- (5,5.25) -- (4.75,5.25) -- cycle}
\fxfcell{4}{5.25}{4.25}{5.5}{1}
\fxfliquid{(4,5.25) -- (4.25,5.25) -- (4.25,5.5) -- (4,5.5) -- cycle}
\fxfcell{4.25}{5.25}{4.5}{5.5}{1}
\fxfliquid{(4.25,5.25) -- (4.5,5.25) -- (4.5,5.5) -- (4.25,5.5) -- cycle}
\fxfcell{4.5}{5.25}{4.75}{5.5}{1}
\fxfliquid{(4.5,5.25) -- (4.75,5.25) -- (4.75,5.5) -- (4.5,5.5) -- cycle}
\fxfcell{4.75}{5.25}{5}{5.5}{1}
\fxfliquid{(4.75,5.25) -- (5,5.25) -- (5,5.5) -- (4.75,5.5) -- cycle}
\fxfcell{4}{5.5}{4.25}{5.75}{1}
\fxfliquid{(4,5.5) -- (4.25,5.5) -- (4.25,5.75) -- (4,5.75) -- cycle}
\fxfcell{4.25}{5.5}{4.5}{5.75}{1}
\fxfliquid{(4.25,5.5) -- (4.5,5.5) -- (4.5,5.75) -- (4.25,5.75) -- cycle}
\fxfcell{4.5}{5.5}{4.75}{5.75}{1}
\fxfliquid{(4.5,5.5) -- (4.75,5.5) -- (4.75,5.75) -- (4.5,5.75) -- cycle}
\fxfcell{4.75}{5.5}{5}{5.75}{1}
\fxfliquid{(4.75,5.5) -- (5,5.5) -- (5,5.75) -- (4.75,5.75) -- cycle}
\fxfcell{4}{5.75}{4.25}{6}{1}
\fxfliquid{(4,5.75) -- (4.25,5.75) -- (4.25,6) -- (4,6) -- cycle}
\fxfcell{4.25}{5.75}{4.5}{6}{1}
\fxfliquid{(4.25,5.75) -- (4.5,5.75) -- (4.5,6) -- (4.25,6) -- cycle}
\fxfcell{4.5}{5.75}{4.75}{6}{1}
\fxfliquid{(4.5,5.75) -- (4.75,5.75) -- (4.75,6) -- (4.5,6) -- cycle}
\fxfcell{4.75}{5.75}{5}{6}{1}
\fxfliquid{(4.75,5.75) -- (5,5.75) -- (5,6) -- (4.75,6) -- cycle}
\fxfcell{5}{5}{5.25}{5.25}{1}
\fxfliquid{(5,5) -- (5.25,5) -- (5.25,5.25) -- (5,5.25) -- cycle}
\fxfcell{5.25}{5}{5.5}{5.25}{1}
\fxfliquid{(5.25,5) -- (5.5,5) -- (5.5,5.25) -- (5.25,5.25) -- cycle}
\fxfcell{5}{5.25}{5.25}{5.5}{1}
\fxfliquid{(5,5.25) -- (5.25,5.25) -- (5.25,5.5) -- (5,5.5) -- cycle}
\fxfcell{5.25}{5.25}{5.5}{5.5}{1}
\fxfliquid{(5.25,5.25) -- (5.5,5.25) -- (5.5,5.5) -- (5.25,5.5) -- cycle}
\fxfcell{5}{5.5}{5.25}{5.75}{1}
\fxfliquid{(5,5.5) -- (5.25,5.5) -- (5.25,5.75) -- (5,5.75) -- cycle}
\fxfcell{5.25}{5.5}{5.5}{5.75}{1}
\fxfliquid{(5.25,5.5) -- (5.5,5.5) -- (5.5,5.75) -- (5.25,5.75) -- cycle}
\fxfcell{5}{5.75}{5.25}{6}{1}
\fxfliquid{(5,5.75) -- (5.25,5.75) -- (5.25,6) -- (5,6) -- cycle}
\fxfcell{5.25}{5.75}{5.5}{6}{1}
\fxfliquid{(5.25,5.75) -- (5.5,5.75) -- (5.5,5.980812) -- (5.465599,6) -- (5.25,6) -- cycle}
\fxfplic{5.5}{5.980812}{5.465599}{6}
\fxfcell{4}{6}{4.25}{6.25}{1}
\fxfliquid{(4,6) -- (4.25,6) -- (4.25,6.25) -- (4,6.25) -- cycle}
\fxfcell{4.25}{6}{4.5}{6.25}{1}
\fxfliquid{(4.25,6) -- (4.5,6) -- (4.5,6.25) -- (4.25,6.25) -- cycle}
\fxfcell{4.5}{6}{4.75}{6.25}{1}
\fxfliquid{(4.5,6) -- (4.75,6) -- (4.75,6.25) -- (4.5,6.25) -- cycle}
\fxfcell{4.75}{6}{5}{6.25}{1}
\fxfliquid{(4.75,6) -- (5,6) -- (5,6.219618) -- (4.887192,6.25) -- (4.75,6.25) -- cycle}
\fxfplic{5}{6.219618}{4.887192}{6.25}
\fxfcell{4}{6.25}{4.25}{6.5}{1}
\fxfliquid{(4,6.25) -- (4.25,6.25) -- (4.25,6.331586) -- (4,6.312301) -- cycle}
\fxfplic{4.25}{6.331586}{4}{6.312301}
\fxfcell{4.25}{6.25}{4.5}{6.5}{1}
\fxfliquid{(4.25,6.25) -- (4.5,6.25) -- (4.5,6.323328) -- (4.25,6.331808) -- cycle}
\fxfplic{4.5}{6.323328}{4.25}{6.331808}
\fxfcell{4.5}{6.25}{4.75}{6.5}{1}
\fxfliquid{(4.5,6.25) -- (4.75,6.25) -- (4.75,6.286416) -- (4.5,6.323551) -- cycle}
\fxfplic{4.75}{6.286416}{4.5}{6.323551}
\fxfcell{4.75}{6.25}{5}{6.5}{1}
\fxfliquid{(4.75,6.25) -- (4.887292,6.25) -- (4.75,6.286976) -- cycle}
\fxfplic{4.887292}{6.25}{4.75}{6.286976}
\fxfcell{5}{6}{5.25}{6.25}{1}
\fxfliquid{(5,6) -- (5.25,6) -- (5.25,6.12001) -- (5,6.220636) -- cycle}
\fxfplic{5.25}{6.12001}{5}{6.220636}
\fxfcell{5.25}{6}{5.5}{6.25}{1}
\fxfliquid{(5.25,6) -- (5.468716,6) -- (5.25,6.121998) -- cycle}
\fxfplic{5.468716}{6}{5.25}{6.121998}
\fxfcell{5}{6.25}{5.25}{6.5}{1}
\fxfcell{5.25}{6.25}{5.5}{6.5}{1}
\fxfjump{1}{1}{2}{4.25}{0.5}{0.134521}{0.397586}{0}
\fxfjump{1}{1}{2}{4.75}{0.5}{0.105154}{0.398146}{0}
\fxfjump{1}{1}{2}{5.25}{0.5}{0.035039}{0.398737}{0}
\fxfjump{1}{1}{2.5}{4.25}{0.5}{0.025533}{0.034506}{0}
\fxfjump{1}{1}{2.5}{4.75}{0.5}{0.055233}{0.082967}{0}
\fxfjump{1}{1}{2.5}{5.25}{0.5}{0.124868}{0.257334}{0}
\fxfjump{1}{1}{2.5}{5.75}{0.5}{0.060092}{0.380372}{0}
\fxfjump{1}{1}{3}{5.25}{0.5}{0.001707}{0.002135}{0}
\fxfjump{1}{1}{3}{5.75}{0.5}{0.090339}{0.16412}{0}
\fxfjump{1}{1}{3}{6.25}{0.5}{0.023173}{0.395913}{0}
\fxfjump{1}{1}{3.5}{5.75}{0.5}{0.003669}{0.004616}{0}
\fxfjump{1}{1}{3.5}{6.25}{0.5}{0.059127}{0.223858}{0}
\fxfjump{1}{1}{4}{5.875}{0.25}{0}{0.022437}{1}
\fxfjump{1}{1}{4}{6.125}{0.25}{0.039114}{0.101579}{0}
\fxfjump{1}{1}{4}{6.375}{0.25}{-0.017527}{-0.006382}{0}
\fxfjump{1}{1}{4.25}{6.375}{0.25}{0.001741}{0.028892}{0}
\fxfjump{1}{1}{4.5}{6.375}{0.25}{-0.007123}{-0.133311}{0}
\fxfjump{1}{1}{4.75}{6.125}{0.25}{-0.002163}{-0.010909}{0}
\fxfjump{1}{1}{4.75}{6.375}{0.25}{-0.014142}{-0.536472}{0}
\fxfjump{1}{1}{5}{6.125}{0.25}{-0.022985}{-0.138253}{0}
\fxfjump{1}{1}{5}{6.375}{0.25}{-0.003203}{-0.748703}{0}
\fxfjump{1}{1}{5.25}{5.875}{0.25}{-0.000417}{-0.00208}{0}
\fxfjump{1}{1}{5.25}{6.125}{0.25}{-0.036929}{-0.410157}{0}
\fxfjump{1}{1}{5.5}{5.875}{0.25}{-0.033923}{-0.216249}{0}
\fxfjump{1}{1}{5.5}{6.125}{0.25}{-0.016857}{-0.779232}{0}
\fxfjump{1}{2}{2.25}{4}{0.5}{0.007279}{0.011069}{0}
\fxfjump{1}{2}{2.25}{4.5}{0.5}{-0.029624}{-0.049238}{0}
\fxfjump{1}{2}{2.25}{5}{0.5}{-0.070613}{-0.201073}{0}
\fxfjump{1}{2}{2.75}{5}{0.5}{-0.001707}{-0.002135}{0}
\fxfjump{1}{2}{2.25}{5.5}{0.5}{-0.035039}{-0.398737}{0}
\fxfjump{1}{2}{2.75}{5.5}{0.5}{-0.092031}{-0.165718}{0}
\fxfjump{1}{2}{3.25}{5.5}{0.5}{-0.003669}{-0.004616}{0}
\fxfjump{1}{2}{5.375}{5.75}{0.25}{-0.000417}{-0.00208}{0}
\fxfjump{1}{2}{2.75}{6}{0.5}{-0.060092}{-0.380372}{0}
\fxfjump{1}{2}{3.25}{6}{0.5}{-0.134374}{-0.297947}{0}
\fxfjump{1}{2}{3.75}{6}{0.5}{-0.078227}{-0.119258}{0}
\fxfjump{1}{2}{4.125}{6}{0.25}{0}{-0.015671}{1}
\fxfjump{1}{2}{4.875}{6}{0.25}{-0.002163}{-0.010909}{0}
\fxfjump{1}{2}{5.125}{6}{0.25}{-0.02515}{-0.14881}{0}
\fxfjump{1}{2}{5.375}{6}{0.25}{-0.061696}{-0.507764}{0}
\fxfjump{1}{2}{4.125}{6.25}{0.25}{-0.054001}{-0.405269}{0}
\fxfjump{1}{2}{4.375}{6.25}{0.25}{-0.054385}{-0.412772}{0}
\fxfjump{1}{2}{4.625}{6.25}{0.25}{-0.061514}{-0.501399}{0}
\fxfjump{1}{2}{4.875}{6.25}{0.25}{-0.073511}{-0.721207}{0}
\fxfjump{1}{2}{5.125}{6.25}{0.25}{-0.053762}{-0.783653}{0}
\fxfjump{1}{2}{5.375}{6.25}{0.25}{-0.016857}{-0.779232}{0}
\fxfjump{1}{2}{3.25}{6.5}{0.5}{-0.023173}{-0.395913}{0}
\fxfjump{1}{2}{3.75}{6.5}{0.5}{-0.082051}{-0.408191}{0}
\fxfjump{1}{2}{4.125}{6.5}{0.25}{-0.021819}{-0.718584}{0}
\fxfjump{1}{2}{4.375}{6.5}{0.25}{-0.024465}{-0.780304}{0}
\fxfjump{1}{2}{4.625}{6.5}{0.25}{-0.017343}{-0.778333}{0}
\fxfjump{1}{2}{4.875}{6.5}{0.25}{-0.003203}{-0.748703}{0}
\fxfjump{2}{1}{2}{4.25}{0.5}{0.160054}{0.214311}{0}
\fxfjump{2}{1}{2}{4.75}{0.5}{0.160387}{0.27512}{0}
\fxfjump{2}{1}{2.5}{5.25}{0.5}{0.161615}{0.15089}{0}
\fxfjump{2}{1}{3}{5.75}{0.5}{0.154197}{0.12691}{0}
\fxfjump{2}{1}{3.5}{6.25}{0.5}{0.160588}{0.275447}{0}
\fxfjump{2}{1}{4}{6.125}{0.25}{0}{-0.035963}{1}
\fxfjump{2}{1}{4}{6.375}{0.25}{-0.078193}{-0.175936}{0}
\fxfjump{2}{1}{5.25}{6.125}{0.25}{-0.078937}{-0.397193}{0}
\fxfjump{2}{2}{2.25}{5}{0.5}{-0.160914}{-0.207587}{0}
\fxfjump{2}{2}{2.75}{5.5}{0.5}{-0.153881}{-0.125533}{0}
\fxfjump{2}{2}{3.25}{6}{0.5}{-0.161216}{-0.161947}{0}
\fxfjump{2}{2}{5.375}{6}{0.25}{-0.07897}{-0.294446}{0}
\fxfjump{2}{2}{4.125}{6.25}{0.25}{-0.075821}{-0.246759}{0}
\fxfjump{2}{2}{4.375}{6.25}{0.25}{-0.078849}{-0.271125}{0}
\fxfjump{2}{2}{4.625}{6.25}{0.25}{-0.078857}{-0.291199}{0}
\fxfjump{2}{2}{4.875}{6.25}{0.25}{-0.078877}{-0.350587}{0}
\fxfjump{2}{2}{5.125}{6.25}{0.25}{-0.078911}{-0.521108}{0}
\fxfjump{2}{2}{3.75}{6.5}{0.5}{-0.160278}{-0.33277}{0}
\fxfjump{2}{2}{4.125}{6.5}{0.25}{0}{-0.067844}{1}
\fxfhanging{4}{4.25}
\fxfhanging{4}{4.75}
\fxfhanging{4}{5.25}
\fxfhanging{4}{5.75}
\fxfhanging{4}{6.25}
%
    \end{scope}
    \fxfnone
    \def\fxfhanging##1##2{\ifdim##2pt<6.5pt\ifdim##2pt>4pt
      \node[tr vertex,minimum size=3.4pt] at (##1,##2) {};\fi\fi}%
    % Generated by docs/techref/figures/generate.jl; do not edit.
\fxfcell{2}{4}{2.5}{4.5}{0}
\fxfliquid{(2.062341,4) -- (2.5,4) -- (2.5,4.5) -- (2.097185,4.5) -- cycle}
\fxfplic{2.062341}{4}{2.097185}{4.5}
\fxfcell{2.5}{4}{3}{4.5}{0}
\fxfliquid{(2.5,4) -- (3,4) -- (3,4.5) -- (2.5,4.5) -- cycle}
\fxfcell{3}{4}{3.5}{4.5}{0}
\fxfliquid{(3,4) -- (3.5,4) -- (3.5,4.5) -- (3,4.5) -- cycle}
\fxfcell{3.5}{4}{4}{4.5}{0}
\fxfliquid{(3.5,4) -- (4,4) -- (4,4.5) -- (3.5,4.5) -- cycle}
\fxfcell{2}{4.5}{2.5}{5}{0}
\fxfliquid{(2.092907,4.5) -- (2.5,4.5) -- (2.5,5) -- (2.251465,5) -- cycle}
\fxfplic{2.092907}{4.5}{2.251465}{5}
\fxfcell{2.5}{4.5}{3}{5}{0}
\fxfliquid{(2.5,4.5) -- (3,4.5) -- (3,5) -- (2.5,5) -- cycle}
\fxfcell{3}{4.5}{3.5}{5}{0}
\fxfliquid{(3,4.5) -- (3.5,4.5) -- (3.5,5) -- (3,5) -- cycle}
\fxfcell{3.5}{4.5}{4}{5}{0}
\fxfliquid{(3.5,4.5) -- (4,4.5) -- (4,5) -- (3.5,5) -- cycle}
\fxfcell{2}{5}{2.5}{5.5}{0}
\fxfliquid{(2.235883,5) -- (2.5,5) -- (2.5,5.410432) -- cycle}
\fxfplic{2.235883}{5}{2.5}{5.410432}
\fxfcell{2.5}{5}{3}{5.5}{0}
\fxfliquid{(2.5,5) -- (3,5) -- (3,5.5) -- (2.558302,5.5) -- (2.5,5.4094) -- cycle}
\fxfplic{2.558302}{5.5}{2.5}{5.4094}
\fxfcell{3}{5}{3.5}{5.5}{0}
\fxfliquid{(3,5) -- (3.5,5) -- (3.5,5.5) -- (3,5.5) -- cycle}
\fxfcell{3.5}{5}{4}{5.5}{0}
\fxfliquid{(3.5,5) -- (4,5) -- (4,5.5) -- (3.5,5.5) -- cycle}
\fxfcell{2}{5.5}{2.5}{6}{0}
\fxfcell{2.5}{5.5}{3}{6}{0}
\fxfliquid{(2.5,5.5) -- (3,5.5) -- (3,5.880728) -- (2.5,5.510641) -- cycle}
\fxfplic{3}{5.880728}{2.5}{5.510641}
\fxfcell{3}{5.5}{3.5}{6}{0}
\fxfliquid{(3,5.5) -- (3.5,5.5) -- (3.5,6) -- (3.141854,6) -- (3,5.919773) -- cycle}
\fxfplic{3.141854}{6}{3}{5.919773}
\fxfcell{3.5}{5.5}{4}{6}{0}
\fxfliquid{(3.5,5.5) -- (4,5.5) -- (4,6) -- (3.5,6) -- cycle}
\fxfcell{2}{6}{2.5}{6.5}{0}
\fxfcell{2.5}{6}{3}{6.5}{0}
\fxfcell{3}{6}{3.5}{6.5}{0}
\fxfliquid{(3.143521,6) -- (3.5,6) -- (3.5,6.201611) -- cycle}
\fxfplic{3.143521}{6}{3.5}{6.201611}
\fxfcell{3.5}{6}{4}{6.5}{0}
\fxfliquid{(3.5,6) -- (4,6) -- (4,6.322031) -- (3.5,6.189898) -- cycle}
\fxfplic{4}{6.322031}{3.5}{6.189898}
\fxfcell{4}{4}{4.25}{4.25}{1}
\fxfliquid{(4,4) -- (4.25,4) -- (4.25,4.25) -- (4,4.25) -- cycle}
\fxfcell{4.25}{4}{4.5}{4.25}{1}
\fxfliquid{(4.25,4) -- (4.5,4) -- (4.5,4.25) -- (4.25,4.25) -- cycle}
\fxfcell{4.5}{4}{4.75}{4.25}{1}
\fxfliquid{(4.5,4) -- (4.75,4) -- (4.75,4.25) -- (4.5,4.25) -- cycle}
\fxfcell{4.75}{4}{5}{4.25}{1}
\fxfliquid{(4.75,4) -- (5,4) -- (5,4.25) -- (4.75,4.25) -- cycle}
\fxfcell{4}{4.25}{4.25}{4.5}{1}
\fxfliquid{(4,4.25) -- (4.25,4.25) -- (4.25,4.5) -- (4,4.5) -- cycle}
\fxfcell{4.25}{4.25}{4.5}{4.5}{1}
\fxfliquid{(4.25,4.25) -- (4.5,4.25) -- (4.5,4.5) -- (4.25,4.5) -- cycle}
\fxfcell{4.5}{4.25}{4.75}{4.5}{1}
\fxfliquid{(4.5,4.25) -- (4.75,4.25) -- (4.75,4.5) -- (4.5,4.5) -- cycle}
\fxfcell{4.75}{4.25}{5}{4.5}{1}
\fxfliquid{(4.75,4.25) -- (5,4.25) -- (5,4.5) -- (4.75,4.5) -- cycle}
\fxfcell{4}{4.5}{4.25}{4.75}{1}
\fxfliquid{(4,4.5) -- (4.25,4.5) -- (4.25,4.75) -- (4,4.75) -- cycle}
\fxfcell{4.25}{4.5}{4.5}{4.75}{1}
\fxfliquid{(4.25,4.5) -- (4.5,4.5) -- (4.5,4.75) -- (4.25,4.75) -- cycle}
\fxfcell{4.5}{4.5}{4.75}{4.75}{1}
\fxfliquid{(4.5,4.5) -- (4.75,4.5) -- (4.75,4.75) -- (4.5,4.75) -- cycle}
\fxfcell{4.75}{4.5}{5}{4.75}{1}
\fxfliquid{(4.75,4.5) -- (5,4.5) -- (5,4.75) -- (4.75,4.75) -- cycle}
\fxfcell{4}{4.75}{4.25}{5}{1}
\fxfliquid{(4,4.75) -- (4.25,4.75) -- (4.25,5) -- (4,5) -- cycle}
\fxfcell{4.25}{4.75}{4.5}{5}{1}
\fxfliquid{(4.25,4.75) -- (4.5,4.75) -- (4.5,5) -- (4.25,5) -- cycle}
\fxfcell{4.5}{4.75}{4.75}{5}{1}
\fxfliquid{(4.5,4.75) -- (4.75,4.75) -- (4.75,5) -- (4.5,5) -- cycle}
\fxfcell{4.75}{4.75}{5}{5}{1}
\fxfliquid{(4.75,4.75) -- (5,4.75) -- (5,5) -- (4.75,5) -- cycle}
\fxfcell{5}{4}{5.25}{4.25}{1}
\fxfliquid{(5,4) -- (5.25,4) -- (5.25,4.25) -- (5,4.25) -- cycle}
\fxfcell{5.25}{4}{5.5}{4.25}{1}
\fxfliquid{(5.25,4) -- (5.5,4) -- (5.5,4.25) -- (5.25,4.25) -- cycle}
\fxfcell{5}{4.25}{5.25}{4.5}{1}
\fxfliquid{(5,4.25) -- (5.25,4.25) -- (5.25,4.5) -- (5,4.5) -- cycle}
\fxfcell{5.25}{4.25}{5.5}{4.5}{1}
\fxfliquid{(5.25,4.25) -- (5.5,4.25) -- (5.5,4.5) -- (5.25,4.5) -- cycle}
\fxfcell{5}{4.5}{5.25}{4.75}{1}
\fxfliquid{(5,4.5) -- (5.25,4.5) -- (5.25,4.75) -- (5,4.75) -- cycle}
\fxfcell{5.25}{4.5}{5.5}{4.75}{1}
\fxfliquid{(5.25,4.5) -- (5.5,4.5) -- (5.5,4.75) -- (5.25,4.75) -- cycle}
\fxfcell{5}{4.75}{5.25}{5}{1}
\fxfliquid{(5,4.75) -- (5.25,4.75) -- (5.25,5) -- (5,5) -- cycle}
\fxfcell{5.25}{4.75}{5.5}{5}{1}
\fxfliquid{(5.25,4.75) -- (5.5,4.75) -- (5.5,5) -- (5.25,5) -- cycle}
\fxfcell{4}{5}{4.25}{5.25}{1}
\fxfliquid{(4,5) -- (4.25,5) -- (4.25,5.25) -- (4,5.25) -- cycle}
\fxfcell{4.25}{5}{4.5}{5.25}{1}
\fxfliquid{(4.25,5) -- (4.5,5) -- (4.5,5.25) -- (4.25,5.25) -- cycle}
\fxfcell{4.5}{5}{4.75}{5.25}{1}
\fxfliquid{(4.5,5) -- (4.75,5) -- (4.75,5.25) -- (4.5,5.25) -- cycle}
\fxfcell{4.75}{5}{5}{5.25}{1}
\fxfliquid{(4.75,5) -- (5,5) -- (5,5.25) -- (4.75,5.25) -- cycle}
\fxfcell{4}{5.25}{4.25}{5.5}{1}
\fxfliquid{(4,5.25) -- (4.25,5.25) -- (4.25,5.5) -- (4,5.5) -- cycle}
\fxfcell{4.25}{5.25}{4.5}{5.5}{1}
\fxfliquid{(4.25,5.25) -- (4.5,5.25) -- (4.5,5.5) -- (4.25,5.5) -- cycle}
\fxfcell{4.5}{5.25}{4.75}{5.5}{1}
\fxfliquid{(4.5,5.25) -- (4.75,5.25) -- (4.75,5.5) -- (4.5,5.5) -- cycle}
\fxfcell{4.75}{5.25}{5}{5.5}{1}
\fxfliquid{(4.75,5.25) -- (5,5.25) -- (5,5.5) -- (4.75,5.5) -- cycle}
\fxfcell{4}{5.5}{4.25}{5.75}{1}
\fxfliquid{(4,5.5) -- (4.25,5.5) -- (4.25,5.75) -- (4,5.75) -- cycle}
\fxfcell{4.25}{5.5}{4.5}{5.75}{1}
\fxfliquid{(4.25,5.5) -- (4.5,5.5) -- (4.5,5.75) -- (4.25,5.75) -- cycle}
\fxfcell{4.5}{5.5}{4.75}{5.75}{1}
\fxfliquid{(4.5,5.5) -- (4.75,5.5) -- (4.75,5.75) -- (4.5,5.75) -- cycle}
\fxfcell{4.75}{5.5}{5}{5.75}{1}
\fxfliquid{(4.75,5.5) -- (5,5.5) -- (5,5.75) -- (4.75,5.75) -- cycle}
\fxfcell{4}{5.75}{4.25}{6}{1}
\fxfliquid{(4,5.75) -- (4.25,5.75) -- (4.25,6) -- (4,6) -- cycle}
\fxfcell{4.25}{5.75}{4.5}{6}{1}
\fxfliquid{(4.25,5.75) -- (4.5,5.75) -- (4.5,6) -- (4.25,6) -- cycle}
\fxfcell{4.5}{5.75}{4.75}{6}{1}
\fxfliquid{(4.5,5.75) -- (4.75,5.75) -- (4.75,6) -- (4.5,6) -- cycle}
\fxfcell{4.75}{5.75}{5}{6}{1}
\fxfliquid{(4.75,5.75) -- (5,5.75) -- (5,6) -- (4.75,6) -- cycle}
\fxfcell{5}{5}{5.25}{5.25}{1}
\fxfliquid{(5,5) -- (5.25,5) -- (5.25,5.25) -- (5,5.25) -- cycle}
\fxfcell{5.25}{5}{5.5}{5.25}{1}
\fxfliquid{(5.25,5) -- (5.5,5) -- (5.5,5.25) -- (5.25,5.25) -- cycle}
\fxfcell{5}{5.25}{5.25}{5.5}{1}
\fxfliquid{(5,5.25) -- (5.25,5.25) -- (5.25,5.5) -- (5,5.5) -- cycle}
\fxfcell{5.25}{5.25}{5.5}{5.5}{1}
\fxfliquid{(5.25,5.25) -- (5.5,5.25) -- (5.5,5.5) -- (5.25,5.5) -- cycle}
\fxfcell{5}{5.5}{5.25}{5.75}{1}
\fxfliquid{(5,5.5) -- (5.25,5.5) -- (5.25,5.75) -- (5,5.75) -- cycle}
\fxfcell{5.25}{5.5}{5.5}{5.75}{1}
\fxfliquid{(5.25,5.5) -- (5.5,5.5) -- (5.5,5.75) -- (5.25,5.75) -- cycle}
\fxfcell{5}{5.75}{5.25}{6}{1}
\fxfliquid{(5,5.75) -- (5.25,5.75) -- (5.25,6) -- (5,6) -- cycle}
\fxfcell{5.25}{5.75}{5.5}{6}{1}
\fxfliquid{(5.25,5.75) -- (5.5,5.75) -- (5.5,5.980812) -- (5.465599,6) -- (5.25,6) -- cycle}
\fxfplic{5.5}{5.980812}{5.465599}{6}
\fxfcell{4}{6}{4.25}{6.25}{1}
\fxfliquid{(4,6) -- (4.25,6) -- (4.25,6.25) -- (4,6.25) -- cycle}
\fxfcell{4.25}{6}{4.5}{6.25}{1}
\fxfliquid{(4.25,6) -- (4.5,6) -- (4.5,6.25) -- (4.25,6.25) -- cycle}
\fxfcell{4.5}{6}{4.75}{6.25}{1}
\fxfliquid{(4.5,6) -- (4.75,6) -- (4.75,6.25) -- (4.5,6.25) -- cycle}
\fxfcell{4.75}{6}{5}{6.25}{1}
\fxfliquid{(4.75,6) -- (5,6) -- (5,6.219618) -- (4.887192,6.25) -- (4.75,6.25) -- cycle}
\fxfplic{5}{6.219618}{4.887192}{6.25}
\fxfcell{4}{6.25}{4.25}{6.5}{1}
\fxfliquid{(4,6.25) -- (4.25,6.25) -- (4.25,6.331586) -- (4,6.312301) -- cycle}
\fxfplic{4.25}{6.331586}{4}{6.312301}
\fxfcell{4.25}{6.25}{4.5}{6.5}{1}
\fxfliquid{(4.25,6.25) -- (4.5,6.25) -- (4.5,6.323328) -- (4.25,6.331808) -- cycle}
\fxfplic{4.5}{6.323328}{4.25}{6.331808}
\fxfcell{4.5}{6.25}{4.75}{6.5}{1}
\fxfliquid{(4.5,6.25) -- (4.75,6.25) -- (4.75,6.286416) -- (4.5,6.323551) -- cycle}
\fxfplic{4.75}{6.286416}{4.5}{6.323551}
\fxfcell{4.75}{6.25}{5}{6.5}{1}
\fxfliquid{(4.75,6.25) -- (4.887292,6.25) -- (4.75,6.286976) -- cycle}
\fxfplic{4.887292}{6.25}{4.75}{6.286976}
\fxfcell{5}{6}{5.25}{6.25}{1}
\fxfliquid{(5,6) -- (5.25,6) -- (5.25,6.12001) -- (5,6.220636) -- cycle}
\fxfplic{5.25}{6.12001}{5}{6.220636}
\fxfcell{5.25}{6}{5.5}{6.25}{1}
\fxfliquid{(5.25,6) -- (5.468716,6) -- (5.25,6.121998) -- cycle}
\fxfplic{5.468716}{6}{5.25}{6.121998}
\fxfcell{5}{6.25}{5.25}{6.5}{1}
\fxfcell{5.25}{6.25}{5.5}{6.5}{1}
\fxfjump{1}{1}{2}{4.25}{0.5}{0.134521}{0.397586}{0}
\fxfjump{1}{1}{2}{4.75}{0.5}{0.105154}{0.398146}{0}
\fxfjump{1}{1}{2}{5.25}{0.5}{0.035039}{0.398737}{0}
\fxfjump{1}{1}{2.5}{4.25}{0.5}{0.025533}{0.034506}{0}
\fxfjump{1}{1}{2.5}{4.75}{0.5}{0.055233}{0.082967}{0}
\fxfjump{1}{1}{2.5}{5.25}{0.5}{0.124868}{0.257334}{0}
\fxfjump{1}{1}{2.5}{5.75}{0.5}{0.060092}{0.380372}{0}
\fxfjump{1}{1}{3}{5.25}{0.5}{0.001707}{0.002135}{0}
\fxfjump{1}{1}{3}{5.75}{0.5}{0.090339}{0.16412}{0}
\fxfjump{1}{1}{3}{6.25}{0.5}{0.023173}{0.395913}{0}
\fxfjump{1}{1}{3.5}{5.75}{0.5}{0.003669}{0.004616}{0}
\fxfjump{1}{1}{3.5}{6.25}{0.5}{0.059127}{0.223858}{0}
\fxfjump{1}{1}{4}{5.875}{0.25}{0}{0.022437}{1}
\fxfjump{1}{1}{4}{6.125}{0.25}{0.039114}{0.101579}{0}
\fxfjump{1}{1}{4}{6.375}{0.25}{-0.017527}{-0.006382}{0}
\fxfjump{1}{1}{4.25}{6.375}{0.25}{0.001741}{0.028892}{0}
\fxfjump{1}{1}{4.5}{6.375}{0.25}{-0.007123}{-0.133311}{0}
\fxfjump{1}{1}{4.75}{6.125}{0.25}{-0.002163}{-0.010909}{0}
\fxfjump{1}{1}{4.75}{6.375}{0.25}{-0.014142}{-0.536472}{0}
\fxfjump{1}{1}{5}{6.125}{0.25}{-0.022985}{-0.138253}{0}
\fxfjump{1}{1}{5}{6.375}{0.25}{-0.003203}{-0.748703}{0}
\fxfjump{1}{1}{5.25}{5.875}{0.25}{-0.000417}{-0.00208}{0}
\fxfjump{1}{1}{5.25}{6.125}{0.25}{-0.036929}{-0.410157}{0}
\fxfjump{1}{1}{5.5}{5.875}{0.25}{-0.033923}{-0.216249}{0}
\fxfjump{1}{1}{5.5}{6.125}{0.25}{-0.016857}{-0.779232}{0}
\fxfjump{1}{2}{2.25}{4}{0.5}{0.007279}{0.011069}{0}
\fxfjump{1}{2}{2.25}{4.5}{0.5}{-0.029624}{-0.049238}{0}
\fxfjump{1}{2}{2.25}{5}{0.5}{-0.070613}{-0.201073}{0}
\fxfjump{1}{2}{2.75}{5}{0.5}{-0.001707}{-0.002135}{0}
\fxfjump{1}{2}{2.25}{5.5}{0.5}{-0.035039}{-0.398737}{0}
\fxfjump{1}{2}{2.75}{5.5}{0.5}{-0.092031}{-0.165718}{0}
\fxfjump{1}{2}{3.25}{5.5}{0.5}{-0.003669}{-0.004616}{0}
\fxfjump{1}{2}{5.375}{5.75}{0.25}{-0.000417}{-0.00208}{0}
\fxfjump{1}{2}{2.75}{6}{0.5}{-0.060092}{-0.380372}{0}
\fxfjump{1}{2}{3.25}{6}{0.5}{-0.134374}{-0.297947}{0}
\fxfjump{1}{2}{3.75}{6}{0.5}{-0.078227}{-0.119258}{0}
\fxfjump{1}{2}{4.125}{6}{0.25}{0}{-0.015671}{1}
\fxfjump{1}{2}{4.875}{6}{0.25}{-0.002163}{-0.010909}{0}
\fxfjump{1}{2}{5.125}{6}{0.25}{-0.02515}{-0.14881}{0}
\fxfjump{1}{2}{5.375}{6}{0.25}{-0.061696}{-0.507764}{0}
\fxfjump{1}{2}{4.125}{6.25}{0.25}{-0.054001}{-0.405269}{0}
\fxfjump{1}{2}{4.375}{6.25}{0.25}{-0.054385}{-0.412772}{0}
\fxfjump{1}{2}{4.625}{6.25}{0.25}{-0.061514}{-0.501399}{0}
\fxfjump{1}{2}{4.875}{6.25}{0.25}{-0.073511}{-0.721207}{0}
\fxfjump{1}{2}{5.125}{6.25}{0.25}{-0.053762}{-0.783653}{0}
\fxfjump{1}{2}{5.375}{6.25}{0.25}{-0.016857}{-0.779232}{0}
\fxfjump{1}{2}{3.25}{6.5}{0.5}{-0.023173}{-0.395913}{0}
\fxfjump{1}{2}{3.75}{6.5}{0.5}{-0.082051}{-0.408191}{0}
\fxfjump{1}{2}{4.125}{6.5}{0.25}{-0.021819}{-0.718584}{0}
\fxfjump{1}{2}{4.375}{6.5}{0.25}{-0.024465}{-0.780304}{0}
\fxfjump{1}{2}{4.625}{6.5}{0.25}{-0.017343}{-0.778333}{0}
\fxfjump{1}{2}{4.875}{6.5}{0.25}{-0.003203}{-0.748703}{0}
\fxfjump{2}{1}{2}{4.25}{0.5}{0.160054}{0.214311}{0}
\fxfjump{2}{1}{2}{4.75}{0.5}{0.160387}{0.27512}{0}
\fxfjump{2}{1}{2.5}{5.25}{0.5}{0.161615}{0.15089}{0}
\fxfjump{2}{1}{3}{5.75}{0.5}{0.154197}{0.12691}{0}
\fxfjump{2}{1}{3.5}{6.25}{0.5}{0.160588}{0.275447}{0}
\fxfjump{2}{1}{4}{6.125}{0.25}{0}{-0.035963}{1}
\fxfjump{2}{1}{4}{6.375}{0.25}{-0.078193}{-0.175936}{0}
\fxfjump{2}{1}{5.25}{6.125}{0.25}{-0.078937}{-0.397193}{0}
\fxfjump{2}{2}{2.25}{5}{0.5}{-0.160914}{-0.207587}{0}
\fxfjump{2}{2}{2.75}{5.5}{0.5}{-0.153881}{-0.125533}{0}
\fxfjump{2}{2}{3.25}{6}{0.5}{-0.161216}{-0.161947}{0}
\fxfjump{2}{2}{5.375}{6}{0.25}{-0.07897}{-0.294446}{0}
\fxfjump{2}{2}{4.125}{6.25}{0.25}{-0.075821}{-0.246759}{0}
\fxfjump{2}{2}{4.375}{6.25}{0.25}{-0.078849}{-0.271125}{0}
\fxfjump{2}{2}{4.625}{6.25}{0.25}{-0.078857}{-0.291199}{0}
\fxfjump{2}{2}{4.875}{6.25}{0.25}{-0.078877}{-0.350587}{0}
\fxfjump{2}{2}{5.125}{6.25}{0.25}{-0.078911}{-0.521108}{0}
\fxfjump{2}{2}{3.75}{6.5}{0.5}{-0.160278}{-0.33277}{0}
\fxfjump{2}{2}{4.125}{6.5}{0.25}{0}{-0.067844}{1}
\fxfhanging{4}{4.25}
\fxfhanging{4}{4.75}
\fxfhanging{4}{5.25}
\fxfhanging{4}{5.75}
\fxfhanging{4}{6.25}
%
    \draw[black,line width=0.6pt] (2.0,4.0) rectangle (5.5,6.5);
    \draw[black!70,dash pattern=on 2pt off 1.5pt,line width=0.5pt] (4,3.9) -- (4,6.6);
  }%
  \begin{subfigure}[b]{0.495\textwidth}
    \centering
    \begin{tikzpicture}[x=1.78cm,y=1.78cm]
      \fxfpanel{1}
      \node[tr small,anchor=north west,fill=white,fill opacity=0.85,text opacity=1] at (2.02,6.48)
        {$\fjump{\surftens}_f=\surftens\curvface_f(\jumpmat\condfrac{\liq}{})_f$};
    \end{tikzpicture}
    \caption{CSF}
    \label{fig:forces:jumps:csf}
  \end{subfigure}\hfill
  \begin{subfigure}[b]{0.495\textwidth}
    \centering
    \begin{tikzpicture}[x=1.78cm,y=1.78cm]
      \fxfpanel{2}
      \node[tr small,anchor=north west,fill=white,fill opacity=0.85,text opacity=1] at (2.02,6.48)
        {$\fjump{\surftens}_f=\surftens\curvface_f(\jumpmat\indicator{\liq})_f$};
    \end{tikzpicture}
    \caption{CGFM}
    \label{fig:forces:jumps:cgfm}
  \end{subfigure}
  \caption{Surface-tension jumps and accelerations on the composite drop
    fixture of \cref{sec:forces:surface-tension} (level transition dashed,
    hanging vertices open circles; liquid shaded). Red: faces with
    $\fjump{\surftens}_f\neq0$~\cref{eq:forces:st-jump}, \chkForceCsfFaces\
    for CSF and \chkForceCgfmFaces\ for CGFM over the whole mesh; arrows: the
    face acceleration $(\forcemap\fjump{\surftens})_f$~\cref{eq:forces:map},
    scaled by its maximum. CSF spreads the jump over the faces between cells
    with different conditional fractions, CGFM puts it on the faces where the
    indicator changes. Purple: faces without a jump whose acceleration is not
    zero because the metric $\rieszmat$ couples the segments of a vertex patch
    at the level transition (\chkForceCsfCoupled\ for CSF, \chkForceCgfmCoupled\
    for CGFM); the pressure gradient couples them in the same way.}
  \label{fig:forces:jumps}
\end{figure}
